% Options for packages loaded elsewhere
\PassOptionsToPackage{unicode}{hyperref}
\PassOptionsToPackage{hyphens}{url}
\documentclass[
]{article}
\usepackage{xcolor}
\usepackage{amsmath,amssymb}
\setcounter{secnumdepth}{-\maxdimen} % remove section numbering
\usepackage{iftex}
\ifPDFTeX
  \usepackage[T1]{fontenc}
  \usepackage[utf8]{inputenc}
  \usepackage{textcomp} % provide euro and other symbols
\else % if luatex or xetex
  \usepackage{unicode-math} % this also loads fontspec
  \defaultfontfeatures{Scale=MatchLowercase}
  \defaultfontfeatures[\rmfamily]{Ligatures=TeX,Scale=1}
\fi
\usepackage{lmodern}
\ifPDFTeX\else
  % xetex/luatex font selection
\fi
% Use upquote if available, for straight quotes in verbatim environments
\IfFileExists{upquote.sty}{\usepackage{upquote}}{}
\IfFileExists{microtype.sty}{% use microtype if available
  \usepackage[]{microtype}
  \UseMicrotypeSet[protrusion]{basicmath} % disable protrusion for tt fonts
}{}
\makeatletter
\@ifundefined{KOMAClassName}{% if non-KOMA class
  \IfFileExists{parskip.sty}{%
    \usepackage{parskip}
  }{% else
    \setlength{\parindent}{0pt}
    \setlength{\parskip}{6pt plus 2pt minus 1pt}}
}{% if KOMA class
  \KOMAoptions{parskip=half}}
\makeatother
\usepackage{longtable,booktabs,array}
\newcounter{none} % for unnumbered tables
\usepackage{calc} % for calculating minipage widths
% Correct order of tables after \paragraph or \subparagraph
\usepackage{etoolbox}
\makeatletter
\patchcmd\longtable{\par}{\if@noskipsec\mbox{}\fi\par}{}{}
\makeatother
% Allow footnotes in longtable head/foot
\IfFileExists{footnotehyper.sty}{\usepackage{footnotehyper}}{\usepackage{footnote}}
\makesavenoteenv{longtable}
\setlength{\emergencystretch}{3em} % prevent overfull lines
\providecommand{\tightlist}{%
  \setlength{\itemsep}{0pt}\setlength{\parskip}{0pt}}
\usepackage{bookmark}
\IfFileExists{xurl.sty}{\usepackage{xurl}}{} % add URL line breaks if available
\urlstyle{same}
\hypersetup{
  hidelinks,
  pdfcreator={LaTeX via pandoc}}

\author{}
\date{}

\begin{document}

\section{Uncertainty Bounding as the Basis of Intelligence: A Formal
Theory}\label{uncertainty-bounding-as-the-basis-of-intelligence-a-formal-theory}

Patrick McCarthy

\begin{center}\rule{0.5\linewidth}{0.5pt}\end{center}

\subsection{Abstract}\label{abstract}

We propose that intelligence is not computational capacity, information
storage, or behavioral breadth, but the \textbf{driven capacity to
improve}: agency \(\mathcal{A}\) --- the intrinsic drive to reduce
uncertainty --- exercised through a higher-order function space \(M\)
that continuously refines both the knowledge graph \(K\) and the
informed action space \(A\). Intelligence is not measured by the size of
what has been learned but by the efficiency with which \(M\) converts
gradient signals into improved \((A, K)\):
\(I(E) = \mathcal{A} \cdot \eta_M\).

Any entity persisting in world state space \(W\) operates under
constraints imposed by \(W\) itself. The active context window \(L_n\)
has a physically bounded capacity \(C_n\); growing it is costly and
eventually impossible. This constraint has a structural consequence that
prior accounts of intelligence do not address: intelligence cannot grow
by filling \(L_n\) --- it must grow by encoding proven knowledge at
lower levels, freeing \(L_n\) for genuine novelty. We show that
graph-structured \(K\) is the most efficient architecture for this.
Graph traversal delivers exactly the knowledge relevant to the current
problem at cost proportional to the problem, not to total \(K\):
progressive disclosure. Selection pressure produces graph-structured
knowledge because it is the most efficient structure under which
\(I(E)\) grows without \(C_n\) becoming a bottleneck.

This gives a formal account of why top-down attention allocation ---
specifically the precision allocation model of the Free Energy Principle
{[}Friston2010, Friston2017{]} --- requires escalation to remain viable
under selection pressure. Top-down allocation has an \(O(k)\) cost in
high-level capacity per concurrent delegation. As \(|A|\) grows,
delegation overhead fills \(C_n\) and inference quality \(\rho\)
degrades --- unless an independent mechanism bounds \(k\). Escalation is
that mechanism: lower-level encoding handles routine problems below
\(L_n\), keeping \(k\) small enough that top-down allocation retains
capacity for genuinely novel problems. \(L_n\) is freed, not burdened,
as the knowledge graph fills. Attention becomes available when not
demanded from below.

From this foundation we derive: the inseparability of knowledge and
action (\(K\) is necessarily indexed by \(A\), with an asymmetric
retention criterion proven from survival pressure), the encoding
hierarchy and its rigor thresholds, multi-timescale gradient descent as
the same process at different learning rates, and the type
irreducibility of \(A\) and \(M\). Extensions to collective
intelligence, sentience, and superintelligence are developed in a
companion paper.

\begin{center}\rule{0.5\linewidth}{0.5pt}\end{center}

\subsection{1. Introduction}\label{introduction}

This framework was not derived by analogy from biology. It emerged from
a concrete engineering problem: generating repeatable, transposable
tasks across a diverse codebase ecosystem spanning thousands of
repositories. The formal theory describes what a self-optimizing agentic
system converged to under real performance constraints. The biological
and civilizational correspondences are predictions the framework makes
--- not the premises it rests on.

\textbf{The origin: provenance as the root constraint.} The initial
problem was repeatability. An agentic system that synthesizes its
behavior from a context window is inherently unpredictable: that
unpredictability is its function --- it exists to handle novelty. But
when you need an action to be transposable across different contexts
(does the recipe that updated service \(A\) with configuration \(X\)
also work for service \(B\) with configuration \(Y\)?), unpredictability
is a liability. Transposability requires a quantifiable basis: not that
the agent believes the action will work, but that there is evidence it
has worked, attributed to a source, grounded in the conditions under
which it succeeded. Without provenance --- without knowing \emph{why}
something worked --- you cannot measure confidence, and without a
confidence measure you cannot optimize. The system cannot learn; it can
only repeat or guess.

This is the engineering form of a general constraint: knowledge transfer
without provenance cannot be reliably weighted, and below some fidelity
threshold the knowledge graph degrades rather than grows. Provenance is
not an epistemic nicety --- it is the minimum infrastructure for a
learning system to function.

\textbf{The graph structure as an optimization.} As the world \(W\) ---
the number of repositories, configurations, and service topologies ---
expanded, the context window problem became the central performance
constraint. Encoding what you know together with how to apply it does
not scale: a node whose knowledge is embedded with its application
context must traverse back to all its callers when any shared element
changes. Shared knowledge cannot be maintained consistently because it
is not represented as shared --- it is duplicated at every point of use.

The knowledge graph solved this by separating \emph{why you know
something} from \emph{how to apply it}. A proposition about what makes
an update strategy work --- why it works, under what conditions ---
becomes a node in \(K\) that multiple informed actions in \(A\) can
reference without embedding. When evidence updates the proposition, all
dependent actions benefit simultaneously. The graph structure was not a
design choice; it was the convergent solution to the problem of
maintaining consistency across a space of actions that grows faster than
any single context can hold.

\textbf{The empirical measurement of factor extraction.} The engineering
consequence of separating \(A\) from \(K\) is precisely measurable.
Before separation, each entity performing a search carried its full
background knowledge implicitly --- \(N\) entities each performing \(M\)
searches yields \(O(N \times M)\) retrieval cost, with no sharing of
common structure. After separation, each entity retrieves its relevant
knowledge by local graph traversal from the action node, bounded by the
problem: \(O(1)\) local in-memory lookups, independent per entity, no
global search, no inter-entity coordination. The
\(O(N \times M) \to O(1)\) transformation is the empirical measurement
of \(K\) precipitating from \(A\) as the normalization of shared
propositions. The 50\% reduction in task-definition code that followed
was a consequence, not the finding; the finding is the scaling law.

\textbf{The escalation architecture as a consequence.} Adding agents
introduced a further optimization that was not available in a purely
knowledge-graph structure: the encapsulation of actions, not just
knowledge. An agent whose scope is a well-defined subtask does not need
access to the full context of the system --- only to the subgraph
\(K^{[a]}\) relevant to its task. This is progressive disclosure in
practice. Subproblems that previously required explicit orchestration
became self-contained agents that escalated to the coordinator only on
failure. As the system matured, proven subnodes converged toward
determinism and were scripted --- promoted to a lower encoding level,
removing them from the space of active inference entirely. Agents that
previously handled uncertain subtasks became deterministic scripts. The
active inference load stayed bounded as the system's coverage of \(W\)
expanded.

\textbf{What this means for the formal claims.} The framework's central
results --- that graph-structured \(K\) is the most efficient
architecture for bounded-context intelligence, that top-down
orchestration degrades inference quality as capability grows, that
knowledge and action are inseparable under a provenance-grounded
retention criterion, and that repeatability (certainty) is the
asymptotic product of learning --- are not theoretical proposals. They
describe what a real system, optimizing against real performance
constraints, independently arrived at. The formal theory is the
explanation of why. The lethality threshold \(U_{lethal}^d\) in this
system is concrete: exceeding token budget, latency threshold, or error
rate across the repository fleet. Selection pressure was cost, not
death. The gradient was real. The convergence was measured.

The test of whether this is a theory of software optimization or a
theory of intelligence is whether the formal structure is consistent
with biology --- the only system where intelligence is unambiguously
instantiated over evolutionary time. We do not claim to prove how
biology works. We claim that if this theory is correct, the biological
architecture should be feasible under it: that genetic encoding of \(M\)
rather than \((K, A)\) is the predicted transmission mode (Section 11,
Open Question 8); that the encoding hierarchy from reflex to active
reasoning follows from the cost structure of Theorem 4 rather than
requiring separate explanation; that cultural transmission --- language,
writing, institutions --- emerges as a structural necessity when the
heritable channel cannot carry individual \((K, A)\) across generations.
None of these were assumed. They follow from the same
optimization constraints that produced the software architecture.
Biology is the feasibility test: if a known intelligent system cannot be
accommodated by the model, the model is wrong. If it can, the model is a
candidate theory of intelligence.

\textbf{On the relationship to prior accounts.} The framework's
relationship to existing theories --- evolutionary biology, the Free
Energy Principle, learning theory --- is neither analogy nor extension.
It is generalization by containment. The framework defines a structural
class: the set of all processes satisfying bounded active context, a
survival threshold, and a growing reachable space. Biological evolution
is one member of that class. Gradient descent in artificial systems is
another. Collective human knowledge transmission is another. This is not
the kind of generalization relativity made of Newtonian mechanics ---
which narrowed Newton by revealing its domain of validity. It is the
kind group theory made of rotational symmetry: defining the abstract
structure that specific instances share, making their common properties
derivable rather than separately observed. The reason evolutionary
systems, learning machines, and collective intelligence all converge on
graph-structured knowledge, escalation architectures, and provenance
requirements is not coincidence --- it is that these are necessary
properties of any member of the class (Theorem 6). The framework does
not explain evolution or AI separately. It explains why they are the
same thing at different timescales and substrates.

\begin{center}\rule{0.5\linewidth}{0.5pt}\end{center}

\subsection{2. The Central Claims}\label{the-central-claims}

We claim: \textbf{intelligence is the driven capacity to improve} ---
agency \(\mathcal{A}\) exercised through a higher-order function space
\(M\) that continuously refines both the knowledge graph \(K\) and the
informed action space \(A\). Formally:
\(I(E) = \mathcal{A} \cdot \eta_M\).

Three components are jointly necessary and none is sufficient alone:

\begin{itemize}
\tightlist
\item
  \textbf{Drive (\(\mathcal{A}\)):} the intrinsic orientation toward
  uncertainty reduction --- the desire to know --- without which \(M\)
  has no reason to engage the gradient
\item
  \textbf{Mechanism (\(\eta_M\)):} the efficiency of \(M\): how
  effectively it converts gradient signals into improved \((A, K)\).
  Without \(\eta_M > 0\), the entity cannot learn from evidence; it
  executes a fixed \((A, K)\) no matter how motivated
\item
  \textbf{Outcome:} the bounding of uncertainty through
  knowledge-grounded action --- what \(I(E)\) produces over time, not
  what it is. The size of \(K\) or \(A\) measures past intelligence;
  \(I(E)\) measures the rate at which that past is being extended
\end{itemize}

\(\mathcal{A} > 0\) and \(\eta_M > 0\) are the only universal invariants
of intelligence (Theorem 6). Every other aspect of \(E\) --- \(A\),
\(K\), substrate, sensing function --- can differ arbitrarily.
\(\mathcal{A} = 0\) or \(\eta_M = 0\) yields the automaton: executing
what prior gradient descent encoded, but not itself descending.

\textbf{The contribution.} Prior theories identify individual aspects of
intelligent systems: FEP derives uncertainty minimization from
self-organization {[}Friston2010{]}; Hutter defines optimal behavior
over computable environments {[}Hutter2005{]}; Chollet measures
skill-acquisition efficiency {[}Chollet2019{]}; Simon establishes
bounded rationality {[}Simon1955{]}. Each takes its constraint as
given. This paper identifies the \emph{generative constraint} --- bounded
active context \(C_n\) under survival pressure in a world with positive
entropy --- and derives the structural consequences as a connected chain:
graph-structured \(K\) (Theorems 2b, 2c), escalation as the necessary
complement to top-down allocation (Theorem 2a), the encoding hierarchy (Theorem 4),
the retention criterion (Theorem 1), and the necessity of
\(\mathcal{A}\) and \(\eta_M\) (Theorem 6). The individual results
formalize known intuitions; the contribution is that they are all
consequences of the same constraint, derived rather than separately
assumed.

\begin{center}\rule{0.5\linewidth}{0.5pt}\end{center}

\subsection{3. Formal Definitions}\label{formal-definitions}

\textbf{Definition 1 (World State Space).} Let \(W\) be a measurable
space of possible world states. At time \(t\), the true state
\(w_t \in W\) is not directly observable by any entity.

\textbf{Definition 2 (Knowledge Graph).} A knowledge graph \(K\) is a
directed graph where nodes are propositions about \(W\) and edges
represent inferential or causal relationships between propositions.
\(K\) is not a passive store; its retention criterion is established by
Theorem 1.

\textbf{Definition 2a (Bounded Knowledge Graph).} A knowledge graph
\(K\) has capacity \(K_{max}\): the maximum number of propositions \(K\)
can hold simultaneously, determined by the substrate's physical
constraints. When \(|K| < K_{max}\), \(K\) operates in the
\textbf{growth regime}: propositions accumulate freely subject to the
retention criterion of Theorem 1. When \(|K| = K_{max}\), \(K\) operates
in the \textbf{curation regime}: adding a new proposition requires
removing an existing one. The two regimes have different retention
criteria.

\textbf{Definition 3 (Informed Action Space).} Let
\(A \subseteq (K \times W_{obs} \to \Delta \Omega)\) be a set of functions,
where \(W_{obs}\) is an entity's observable projection of \(W\) and
\(\Delta \Omega\) is a distribution over possible actions. Each element
\(a \in A\) has type \(f : K \times W_{obs} \to \Delta \Omega\). \(A\) is the
complete set of informed actions available to an entity. An
\textbf{informed action} \(a \in A\) is an action grounded in \(K\),
validated against \(W\), and attributed to its origins: one that knows
why it works. An uninformed action (grounded in no \(K\)) is not an
element of \(A\). Knowledge with no action potential is not an element
of \(A\) either; it is data. An informed action is specifically the
coupling: knowledge that enables action, action that is justified by
knowledge.

\(A\) carries graph structure. The elements of \(A\) are nodes; typed
relations between them are edges. This structure is not a design choice
--- it is forced by the same constraints (retrieval
\(O(|A^{[a]}|)\), bounded \(C_n\), deduplication) that force \(K\) to be
a directed graph (Theorem 2b Part 3), and independently by the
requirement for sub-linear retention evaluation with cascade (Theorem
2c). The proof is identical: any structure satisfying all constraints
simultaneously is a directed graph. Edge
types on \(A\) follow from the same evidence-grounded logic as edge
types on \(K\): hard dependencies between actions (feasibility
constraints), probabilistic ordering evidence (soft gradient
directions), and joint-execution synergies (combined actions that
outperform either alone).

\(A\) as a graph is discrete (nodes and typed edges); each individual
\(a \in A\) maps into the continuous probability simplex \(\Delta \Omega\).
Gradient descent in \(A\) is traversal on the discrete graph, directed
by the continuous-valued \(G(f', K)\) at each adjacent node. \(K\) and
\(A\) are both directed graphs under bounded context and selection
pressure, differing only in domain: \(K\) is what the entity knows;
\(A\) is what it can do.

In general, \(K\) and \(A\) are distinct graphs with many-to-many
coupling. In the degenerate case where every proposition is directly
actionable, \(A = K\). This is the limit of crystallization at the
bottom of the encoding hierarchy: a reflex is simultaneously knowledge
and action. \(A = K\) holds locally in fully crystallized domains;
\(A \neq K\) persists globally because \(\mathcal{T}_{reachable}\) grows
without bound (Theorem 3, Part II).

\textbf{Definition 3a (Action Feedback).} For any \(a \in A\) executed
in world state \(w\), the feedback \(\phi(f, w) \in \Phi\) is the
observable consequence of \(f\) in \(w\): the signal \(W\) returns in
response to the action. The updated knowledge after executing \(f\) and
observing its feedback is:

\[K_f = K \cup \{p_f\}\]

where \(p_f\) is the proposition grounded by observing \(\phi(f, w)\).
The feedback \(\phi(f,w)\) is the formal object underlying the gradient
signal \(\delta\) in Theorem 3: \(\delta =
\phi(f, w) - \mathbb{E}[\phi(f, \cdot) \mid K]\).

\textbf{Definition 3b (Information Gain).} The information gain of
action \(f\) given current knowledge \(K\) is:

\[G(f, K) = \mathbb{E}_w\!\left[U(w, K) - U(w, K_f)\right] = H(W \mid K) - H(W \mid K,\, \phi(f, w)) = I(f\,;\, W \mid K)\]

\(G(f, K) \geq 0\) by the data processing inequality
{[}Cover2006{]}: conditioning on true information cannot increase
conditional entropy. \(G(f, K) = 0\) when \(f\) provides no new
information about \(w\) given \(K\). Corrupted feedback may increase \(U\) --- a proposition grounded in
false evidence degrades \(K\) rather than improving it. Recovery from
corrupted feedback is handled not by the convergence mechanism (Theorem
3, which assumes true information) but by the maintenance mechanism
(Theorem 2c): stored typed adjacency enables the entity to identify
propositions grounded through the corrupted evidence, cascade-remove or
re-evaluate their dependents, and scope the damage to the actual
dependency set rather than the full graph. The two mechanisms are
complementary: Theorem 3 establishes convergence under clean signal;
Theorem 2c provides the corrective architecture when signals are dirty.

\textbf{Definition 4 (Entity).} An entity
\(E = (K, A, \sigma, \mathcal{A})\) where
\(\sigma : W \rightarrow W_{obs}\) is a sensing function projecting the
true world state into the entity's observable space. \(K\) and \(A\)
evolve over time via \(M\); their initial values \((K_0, A_0)\) are
inherited from \(L_0\) encoding. The bootstrapping problem is resolved
by \(L_0\): the entity begins from the encoded product of gradients that
preceded it. Throughout this paper \(K\) and \(A\) refer to
current-time values; the subscript \(0\) denotes initial conditions.

\textbf{Definition 5 (Uncertainty).} The uncertainty of entity \(E\)
about world state \(w\) is the surprisal of \(w\) given \(K\)
{[}Shannon1948, Cover2006{]}:

\[U(w, K) = -\log P(w \mid K)\]

\(K\) is a directed graph of propositions interpreted as a Bayesian
network {[}Pearl1988{]}. Each proposition \(p \in K\) contributes a
likelihood factor \(\ell_p : W \to [0,1]\). The graph's edge structure
encodes conditional independence: propositions with no path between them
are conditionally independent given shared ancestors. The posterior over
world states given the full knowledge graph is:

\[P(w \mid K) \propto P_0(w) \cdot \prod_{p \in K} \ell_p(w)\]

This is the standard Bayes posterior obtained by conditioning \(P_0\) on
all propositions simultaneously. \(P(w \mid K)\) is well-defined and
normalizable for any proper prior \(P_0\) and bounded likelihood
functions \(\ell_p\). The theory is parametric in \(P_0\) --- it does
not depend on a specific prior, only on its existence and
well-formedness.

The expected uncertainty over all world states is the conditional
entropy:

\[H(W \mid K) = \mathbb{E}_w[U(w, K)] = -\sum_w P(w \mid K) \log P(w \mid K)\]

The minimization imperative (Theorem 1) operates on this expected
quantity; \(U(w, K)\) is the pointwise uncertainty the entity faces when
it encounters specific world state \(w\).

<<<<<<< HEAD
\textbf{Definition 6 (Agency).} Agency \(\mathcal{A} \in [0, 1]\) is a
continuous scalar measuring the intrinsic drive of an entity to reduce
its own uncertainty: the desire to know {[}Schmidhuber2010{]}. It is not
an external objective imposed on \(E\), nor a property derived from
\(A\) or \(K\). It is constitutive: \(\mathcal{A}\) is what makes \(A\)
an active informed action space rather than a static structure. At
\(\mathcal{A} = 0\), \(F\) is a set of possible functions with no
internal reason to engage \(W\). At \(\mathcal{A} = 1\), \(F\) is
maximally oriented toward the unbounded frontier. Intermediate values
are meaningful: the drive can be stronger or weaker, shaped by
=======
\textbf{Definition 5a (Survival Threshold).} For entity \(E\) in domain
\(d \subseteq W\), there exists a threshold \(U_{lethal}^d\) such that:

\[U(w, K) > U_{lethal}^d \Rightarrow E \text{ does not survive in } d\]

Above this threshold, uncertainty is so high that the entity cannot act
effectively enough to persist. \(U_{lethal}^d\) is not a choice or a
preference; it is imposed by \(W\).

\textbf{Definition 6 (Agency).} Agency \(\mathcal{A} \in [0, 1]\) is the
probability that entity \(E\) engages its learning mechanism \(M\) in
response to a gradient signal from \(W\): the response gate. It is not
an external objective imposed on \(E\), nor a property derived from
\(A\) or \(K\). It is constitutive: \(\mathcal{A}\) is what makes \(A\)
an active informed action space rather than a static structure. At
\(\mathcal{A} = 0\), \(A\) is a set of possible functions with no
internal reason to engage \(W\). At \(\mathcal{A} = 1\), \(E\) engages
on every gradient signal it encounters. Intermediate values are
meaningful: the response probability can be higher or lower, shaped by
>>>>>>> 319f447daa46f8c1e949fb7929535d1aa6396025
architecture, encoding history, and the density of unresolved
uncertainty in the entity's reachable \(W\).

<<<<<<< HEAD
Agency is the reason gradient descent has a runner. The landscape of
uncertainty exists independently of any entity; \(\mathcal{A}\) is what
determines how strongly an entity moves through it.

\(\mathcal{A}\) is operationalizable: it is proportional to the rate of
uncertainty reduction as survival pressure approaches zero --- that is,
as \(U_{lethal}^d \to \infty\) (the threshold above which death occurs
goes to infinity, meaning no level of uncertainty is lethal):
=======
\(\mathcal{A}\) has an empirical operationalization:
>>>>>>> 319f447daa46f8c1e949fb7929535d1aa6396025

\[\mathcal{A}(E) \propto \lim_{U_{lethal}^d \to \infty} \frac{dU}{dt}(E)\]

An entity with genuine \(\mathcal{A} > 0\) continues engaging \(M\)
when survival pressure is removed; an entity with \(\mathcal{A} = 0\)
stops. This isolates intrinsic engagement from survival-compelled
updating and is measurable in any system where survival pressure can be
varied independently of capacity.

\(\mathcal{A}\) precipitates from \(M\) by the same normalization
argument that precipitates \(K\) from \(A\). An entity that embeds its
directional drive in every application of \(M\) recomputes orientation
at every step --- the \(M\)-analog of the \(O(n \cdot k)\) scaling
failure. An entity that maintains \(\mathcal{A}\) as the shared
directional structure across all \(M\)-applications achieves \(O(1)\)
access to its drive. \(\mathcal{A}\) is therefore the first factor to
precipitate from \(M\)'s cross product: the normalized orientation that
every \(M\)-application shares, extracted for the same reason \(K\) is
extracted from \(A\).

\textbf{Definition 7 (Intelligence).} The intelligence of \(E\) is the
\emph{driven capacity to improve}: the exercise of \(\mathcal{A}\)
through \(M\) to refine \(A\) and \(K\) in response to gradient signals
{[}Schmidhuber2010, Thrun1998{]}:

\[I(E) = \mathcal{A} \cdot \eta_M\]

where \(\mathcal{A} \in [0, 1]\) is the intrinsic drive (Definition 6)
and \(\eta_M\) is the efficiency of the higher-order function space
\(M\) (Definition 11a): the rate at which \(M\) improves \(A\)'s
uncertainty-reduction capacity per unit of learning signal. This
expression is stated here as a definition and established as necessary
in Theorem 6: \(\mathcal{A}\) and \(\eta_M\) are the unique invariants
of persistence under bounded context and selection pressure, and their
product is the minimal measure that captures both.

At \(\mathcal{A} = 0\): \(I(E) = 0\) --- a capable but inert learning
mechanism is not intelligence. At \(\eta_M = 0\): \(I(E) = 0\) --- a
motivated entity whose \(M\) cannot improve \(A\) from evidence is an
automaton.

The expected uncertainty reduction along a trajectory is a
\textbf{performance measure} of current \(A\) and \(K\) --- the
consequence of past \(M\) activity, not the definition of intelligence.
Two entities starting from identical \(K\)
and \(A\) but different \(I(E)\) will diverge: the entity with higher
\(\eta_M\) builds better \(A\) faster, and eventually dominates in
\(U\)-reduction even if initially equal in performance.

Intelligence is not the size of \(K\) or \(A\). \(K\) and \(A\) are what
intelligence has built. Intelligence is the efficiency with which \(M\)
continues building.

<<<<<<< HEAD
=======
\emph{Note on named components.} The components \(K\), \(A\), \(M\), and
\(\mathcal{A}\) named here are dimensions that have precipitated from
the entity's full cross-product state --- call it \(K^n\) --- at the
individual learning level. The theory's general result is that
dimensional precipitation must occur under the joint constraints: as
\(K^n\) grows, shared structure must be extracted and named or retrieval
cost grows without bound. What dimensions precipitate is a function of
the entity's \(W\) and its history; the theory constrains the process,
not the outcome. \(K\), \(A\), \(M\), and \(\mathcal{A}\) are the
dimensions that fall out in entities with active individual learning
mechanisms --- the class this paper studies. Other projections
precipitate differently: biological \(L_0\) encoding (DNA) does not
separate into \(K\), \(A\), \(M\), \(\mathcal{A}\); it is a dense
cross-product that carries all survival-relevant structure without
naming any of it as a distinct component. The named components are an
epistemological convenience grounded in the studied entity class, not a
universal constraint of the theory.

>>>>>>> 319f447daa46f8c1e949fb7929535d1aa6396025
\begin{center}\rule{0.5\linewidth}{0.5pt}\end{center}

\subsection{4. The Inseparability of Knowledge and
Action}\label{the-inseparability-of-knowledge-and-action}

\textbf{Lemma 1 (Proposition Absence Lethality).} For any entity \(E\)
and proposition \(p \in K\) that grounds some \(a \in A\) in domain
\(d \subseteq W\): removing \(p\) from \(K\) weakly increases
uncertainty in \(d\), and in the limiting case where \(p\) is the sole
grounding for the critical action in \(d\), the increase may cause
\(U(w, K \setminus \{p\}) > U_{lethal}^d\), at which point \(E\) does
not survive in \(d\) (Definition 5a).

\emph{Proof.} Conditional entropy is monotonically non-increasing in the
conditioning set {[}Cover2006{]}:
\(H(W \mid K) \leq H(W \mid K \setminus \{p\})\) for any \(p \in K\).
Removing \(p\) weakly increases expected uncertainty. If \(p\) grounds
the only \(a \in A\) applicable in domain \(d\), then
\(K \setminus \{p\}\) contains no action grounded in \(d\): \(A\) has no
mechanism to reduce \(U(w, K \setminus \{p\})\) below its current level
in that domain when \(w \in d\) is encountered. By Definition 5a, if
\(U(w, K \setminus \{p\}) > U_{lethal}^d\) at the encountered \(w\),
\(E\) does not survive. Therefore the cost of absent \(p\) when needed,
\(C_{absent}\), is bounded below by entity non-survival in critical
domains --- which dominates all finite costs. \(\square\)

\textbf{Theorem 1 (K/A Inseparability).} For any entity
\(E = (K, A, \sigma)\), the retention criterion for any proposition
\(p\) is asymmetric:

\begin{itemize}
\tightlist
\item
  \(p\) is \textbf{retained} in \(K\) when \(\exists\, a \in A\)
  applicable to \(p\), or when the survival benefit of \(p\) is
  uncertain
\item
  \(p\) is \textbf{pruned} from \(K\) only when it is certain that \(p\)
  contributes zero to \(P(E, \tau)\) for any \(\tau\), i.e., when it is
  known that no \(a \in A\) can be grounded in \(p\), now or in any
  reachable region of \(W\)
\item
  \(p\) is \textbf{displaced} in the curation regime (\(|K| = K_{max}\))
  when a candidate proposition \(q\) satisfies
  \(\mathbb{E}[G(a_q, K)] > \mathbb{E}[G(a_p, K)]\), where \(p\) is the
  proposition with minimum marginal value in \(K\). Displacement
  requires only that \(q\) contributes more at the margin than \(p\) ---
  not that \(p\) is useless. The asymmetric retention criterion governs
  which propositions enter; the comparative criterion governs which are
  displaced when the graph is full. Information gain \(G\) is submodular
  --- the marginal value of \(p\) depends on what else is in \(K\), so
  the greedy displacement criterion is evaluated at the margin given the
  current full \(K\), not in isolation. This is the correct procedure:
  the entity estimates value relative to what it already holds, not
  relative to an empty graph {[}Nemhauser1978{]}.
\end{itemize}

The criterion is not symmetric. Uncertainty of benefit is not the same
as certainty of no benefit. An entity with \(\mathcal{A}\), whose
\(\mathcal{T}_{reachable}\) grows with every action taken, retains \(p\)
whose value it cannot yet see, because the gradient may reach that
territory. Pruning requires the same standard as certainty (Definition
10): the benefit must be known to be zero across sufficient variation in
\(W\).

\emph{Proof.} Let \(B(p, \tau)\) denote the contribution of \(p\) to
\(P(E, \tau)\) along trajectory \(\tau\), and let
\(\mathcal{T}_{reachable}\) be the set of trajectories reachable from
the entity's current state. Define the option value of retaining \(p\):

\[OV(p) = P\!\left(\exists\,\tau^* \in \mathcal{T}_{reachable} : B(p, \tau^*) > 0\right) \cdot \mathbb{E}\!\left[B(p,\tau^*) \mid B(p, \tau^*) > 0\right]\]

Let \(c_{retain}(p, K) > 0\) denote the total maintenance cost of
holding \(p\) in \(K\). This cost is not merely storage. It includes
three components that grow with the structure of \(K\):

\begin{itemize}
\tightlist
\item
  \textbf{Retrieval interference}: every proposition in \(K\) that is
  not relevant to the current problem contributes noise to the retrieval
  environment. Under graph-structured \(K\), interference is bounded by
  the local degree of \(p\) (only adjacent propositions interact during
  traversal); without graph structure, interference grows with \(|K|\)
  (Theorem 2a).
\item
  \textbf{Maintenance overhead}: every proposition in \(K\) must be
  re-evaluated when evidence changes (Theorem 2c). Holding \(p\) commits
  the entity to evaluating \(p\) during every cascade triggered by
  revisions to propositions in \(p\)'s dependency neighborhood.
\item
  \textbf{Edge maintenance}: retaining \(p\) requires maintaining all of
  \(p\)'s typed edges. Stale or spurious edges degrade cascade precision,
  causing over-retrieval (\(\rho < 1\)) or missed dependencies during
  removal.
\end{itemize}

Let \(c_{prune} \geq 0\) denote the cost of pruning \(p\) (processing,
reorganization). The expected net value of retention is
\(NV_{retain}(p) = OV(p) - c_{retain}(p, K)\).

The expected net value of pruning is:

\[NV_{prune}(p) = -c_{prune} - P\!\left(\exists\,\tau^* : B(p,\tau^*) > 0\right) \cdot C_{absent}\]

where \(C_{absent}\) is the cost incurred when \(p\) is needed but
absent. By Lemma 1, \(C_{absent}\) is bounded below by entity
non-survival in critical domains.

Retention is selected for when \(NV_{retain}(p) > NV_{prune}(p)\). For
propositions that ground critical actions in reachable domains,
\(C_{absent}\) dominates \(c_{retain}\) regardless of maintenance cost,
and retention holds whenever

\[P\!\left(\exists\,\tau^* \in \mathcal{T}_{reachable} : B(p,\tau^*) > 0\right) > 0\]

i.e., whenever survival benefit is uncertain. For non-critical
propositions, however, \(c_{retain}(p, K)\) is not negligible relative
to the expected benefit. As \(|K|\) grows, the maintenance, interference,
and edge costs of marginal propositions increase, and the retention
threshold rises: propositions with low option value are pruned not
because their benefit is certainly zero, but because their maintenance
cost exceeds their marginal expected value. This predicts aggressive
forgetting as correct behavior under bounded capacity --- the entity
prunes low-value propositions to reduce retrieval interference and
maintenance load on the remaining graph, improving \(\rho\) for the
propositions it keeps.

The asymmetry remains for critical propositions: uncertain benefit of a
survival-relevant proposition is sufficient for retention, because
\(C_{absent}\) dominates any finite \(c_{retain}\). But the asymmetry
weakens as propositions move further from the survival boundary, and
vanishes entirely when \(c_{retain}(p, K) > OV(p)\). Forgetting is the
retention criterion working correctly at the margin.

In the curation regime (\(|K| = K_{max}\)), adding \(q\) requires
displacing \(p^* = \arg\min_{p \in K} \mathbb{E}[G(a_p, K)]\).
Displacement occurs when
\(\mathbb{E}[G(a_q, K)] > \mathbb{E}[G(a_{p^*}, K)]\): the asymmetric
criterion applied to a choice between two propositions rather than
between one and nothing. By Theorem 3 Part II,
\(\mathcal{T}_{reachable}\) grows monotonically, keeping future needs
uncertain unless benefit is certainly zero --- preserving the asymmetry
across the entity's lifetime. \(\square\)

\emph{Correspondence.} This formalizes Gibson's affordances
{[}Gibson1979{]} and extends the information bottleneck
{[}Tishby2011{]}, which treats the retention threshold as a design
parameter; here it is derived from survival pressure. The asymmetric
criterion resolves an apparent asymmetry:
an entity retains \(p\) beyond the certainty horizon because the
benefit is uncertain, not known to be zero.

\textbf{Corollary 1 (Indexed Knowledge).} \(K\) is constitutively
indexed by \(A\). Two entities with identical \(\sigma\) but different
\(A\) maintain different \(K\) even when perceiving identical world
states.

\textbf{Corollary 2 (Compactness).} \(K\) is bounded in size by \(A\):
only propositions applicable to available informed actions are retained.

\textbf{Corollary 2a (Curation and Forgetting).} Under bounded \(K\),
forgetting is correct behavior: the curation regime displaces
lower-value propositions as \(\mathcal{T}_{reachable}\) evolves,
ensuring \(K\) tracks the current frontier.

\textbf{Corollary 2b (K as the Normalization of A --- the Scaling
Argument).} An entity can maintain informed actions without a separately
structured \(K\), embedding all propositional knowledge implicitly in
the parameters of each \(a \in A\). This is a valid but scaling-limited
architecture. Let \(n = |A|\) and \(k\) = average number of shared
propositions per action. Without \(K\): context per action is \(O(k)\)
(each \(a\) carries its full background); update cost when one
proposition changes is \(O(n)\) (every \(a\) encoding it must be
individually revised); total cost scales as \(O(n \cdot k)\). With \(K\)
normalized out: context per action is \(O(|K^{[a]}|)\) --- the relevant
subgraph retrieved by traversal (Theorem 2b); update cost per
proposition is \(O(1)\) --- one update to \(K\) propagates to all \(f\)
that reference it; total cost scales as \(O(n + k)\).

The difference \(O(n \cdot k)\) versus \(O(n + k)\) is not a constant
factor --- it is the difference between a system whose context cost
grows multiplicatively with shared knowledge complexity and one whose
context cost stays bounded under growth of \(A\). Under bounded \(C_n\),
the implicit-\(K\) architecture does not scale: every increase in
\(|A|\) increases required context per action proportionally to the
shared background knowledge it carries. \(K\) precipitates from \(A\) as
the normalization of shared propositions --- the factoring-out of common
structure across functions --- because this is the only architecture
that keeps \(C_n\) bounded as \(A\) grows. An entity without \(K\) is
not wrong; it is limited to the regime where \(n \cdot k\) fits in
\(C_n\). Selection pressure at scale removes such entities in favor of
those whose \(K\) has crystallized --- literally, the precipitation of
shared knowledge out of the action space into its own normalized graph.
This is the formal explanation of the \(O(N \times M) \to O(1)\)
empirical measurement described in Section 1: each entity's local graph
traversal is \(O(1)\) precisely because shared propositions have been
factored out of \(A\) into \(K\).

\begin{center}\rule{0.5\linewidth}{0.5pt}\end{center}

\subsection{5. The Escalation Principle}\label{the-escalation-principle}

\textbf{Definition 8 (Encoding Hierarchy).} The encoding hierarchy is a
continuous two-dimensional space \(\mathcal{E}(c, r)\) parameterized by
runtime cost \(c \geq 0\) and reversibility \(r \in [0, 1]\), where
increasing \(c\) corresponds to more expensive engagement and increasing
\(r\) to more easily revised encodings. Consistent with the continuity
principle: a continuous being moves through this space continuously;
there is no categorical jump between encoding modes, only a gradient of
cost and permanence {[}Baars1988, Dehaene2011, Kahneman2011{]}.

Boundaries exist where the qualitative character of encoding changes,
requiring information to be translated as it crosses. These boundaries
are \textbf{interfaces}: formal structures through which signals pass
between qualitatively distinct encoding regimes. The named levels
\(L_0, L_1, \ldots, L_n\) are the interfaces at the significant
boundaries of \(\mathcal{E}(c, r)\):

{\def\LTcaptype{none} % do not increment counter
\begin{longtable}[]{@{}
  >{\raggedright\arraybackslash}p{(\linewidth - 6\tabcolsep) * \real{0.2500}}
  >{\raggedright\arraybackslash}p{(\linewidth - 6\tabcolsep) * \real{0.2500}}
  >{\raggedright\arraybackslash}p{(\linewidth - 6\tabcolsep) * \real{0.2500}}
  >{\raggedright\arraybackslash}p{(\linewidth - 6\tabcolsep) * \real{0.2500}}@{}}
\toprule\noalign{}
\begin{minipage}[b]{\linewidth}\raggedright
Interface
\end{minipage} & \begin{minipage}[b]{\linewidth}\raggedright
Boundary crossed
\end{minipage} & \begin{minipage}[b]{\linewidth}\raggedright
Runtime Cost
\end{minipage} & \begin{minipage}[b]{\linewidth}\raggedright
Reversibility
\end{minipage} \\
\midrule\noalign{}
\endhead
\bottomrule\noalign{}
\endlastfoot
\(L_n\) & Active reasoning / novel uncertainty & \(O(\text{compute})\) &
Fully reversible \\
\(\vdots\) & Learned patterns, habits & Decreasing & Decreasing \\
\(L_1\) & Reflex / autonomic & \(\approx 0\) & Semi-permanent \\
\(L_0\) & Heritable encoding & \(0\) & Permanent within lineage \\
\end{longtable}
}

Escalation (Theorem 2) is the \textbf{upward interface protocol}: the
specification of what form a failure signal takes when crossing a
boundary from lower \(c\) to higher \(c\). Promotion (Theorem 4) is the
\textbf{downward interface protocol}: the condition under which an
action may cross a boundary from higher \(c\) to lower \(c\). The cost
\(C_i\) in Theorem 4 is a property of the interface, not of the
continuous space between interfaces: it is the cost of failure when an
encoding error has crossed the boundary and propagates at the lower
level's permanence.

The number of interfaces is substrate-determined --- how many
qualitative boundaries exist in the entity's physical
\(\mathcal{E}(c, r)\). The theory specifies the structural properties
every interface must have (monotone cost, monotone rigor threshold,
promotion criterion); it does not prescribe how many interfaces exist or
what they are called. The named levels in the table above are empirical
illustrations, not definitional. The continuous space between interfaces
admits arbitrary encodings; the interfaces themselves are discrete
because a boundary is a specific location, not a region.

\textbf{Theorem 2 (Escalation).} An informed action \(a \in A\) is
handled at the lowest encoding level \(L_i\) such that
\(U(w, K^{[a]}) \approx 0\) for the domain of \(w\) relevant to \(f\),
where \(K^{[a]} \subseteq K\) is the knowledge subset applicable to
\(f\). Escalation to \(L_{i+1}\) occurs if and only if \(L_i\) fails to
bound uncertainty.

\emph{Proof.} By Definition 8, the cost of engaging level \(L_i\)
exceeds the cost of engaging \(L_{i-1}\): the encoding hierarchy is
parameterized by increasing runtime cost \(c\), so \(C_{i+1} > C_i\) for
all \(i\). Selection pressure therefore favors handling \(f\) at the
lowest level sufficient to bound uncertainty: engaging \(L_{i+1}\) when
\(L_i\) suffices incurs unnecessary cost without reducing \(U\) further.
An entity that routinely over-escalates incurs higher encoding costs
without survival benefit and is selected against. Conversely, if \(L_i\)
fails to bound \(U(w, K^{[a]})\) --- that is, \(U(w, K^{[a]}) > 0\)
beyond acceptable residual --- then \(f\) is not handled at \(L_i\): the
failure propagates upward until some \(L_j\) (\(j > i\)) contains the
uncertainty or the entity is harmed (Definition 5a). The active context
window \(L_n\) is the level of last resort by construction: it is the
highest-cost, fully reversible level. \(\square\)

\textbf{Corollary 3 (Context Window as Frontier).} \(L_n\) is the
frontier of unbounded uncertainty, not the seat of intelligence. Its
engagement signals failure in lower-level encodings.

\textbf{Corollary 4 (Simultaneous Multi-Level Activity).} All levels
operate concurrently: \(L_n\) handles novel uncertainty while lower
levels handle their domains independently.

\subsubsection{Why the Graph: Progressive Disclosure and the Bounded
Context
Theorem}\label{why-the-graph-progressive-disclosure-and-the-bounded-context-theorem}

\textbf{Definition 8a (Active Context Capacity and Inference Quality).}
Let \(C_n\) denote the capacity of \(L_n\): the maximum number of
propositions simultaneously active in context. \(C_n\) is bounded above
by the substrate's physical constraints --- but the effective capacity
available for any given problem is also shaped by the structure of
\(K\). Graph-structured knowledge enables chunking: a single graph node
can represent a cluster of propositions that co-activate as a unit,
making multiple propositions available at the cost of one context slot.
Expertise increases effective \(C_n\) within a domain by compressing
frequently co-occurring propositions into retrievable chunks --- not by
expanding the substrate bound, but by reducing the number of slots
required per unit of relevant knowledge. The substrate sets the ceiling;
the structure of \(K\) determines how much of \(W\) that ceiling covers.
Let \(\text{ctx}(L_n, t)\) be the set of propositions loaded into
\(L_n\) at time \(t\), with \(|\text{ctx}(L_n, t)| \leq C_n\).

The \textbf{inference quality} of \(L_n\) for problem \(f\) is the
relevance density of its content:

\[\rho(L_n, f) = \frac{|K^{[a]}|}{|\text{ctx}(L_n, t)|}\]

\(\rho = 1\): \(L_n\) holds exactly what \(f\) requires --- maximum
signal, zero noise. \(\rho \to 0\): \(L_n\) is dominated by content
irrelevant to \(f\) --- inference degrades. Every proposition in
\(\text{ctx}(L_n)\) that is not in \(K^{[a]}\) is noise: it contributes
variance to \(L_n\)'s inference without contributing signal. Minimizing
that noise is equivalent to minimizing
\(|\text{ctx}(L_n)| - |K^{[a]}|\).

\textbf{Definition 8b (Delegation Overhead).} Under a top-down
allocation architecture, \(L_n\) directs which subproblems are handled
at lower levels. For each active delegation, \(L_n\) must hold: the
subproblem specification, the expected return signal, and sufficient
context to integrate the result when it arrives. Let \(\alpha > 0\)
denote the minimum context cost per active delegation --- the overhead
top-down allocation imposes on \(L_n\) per managed subproblem,
regardless of whether that subproblem is routine or novel.

\textbf{Theorem 2a (Top-Down Allocation Cost Under Bounded Context).}
Under a top-down allocation architecture with \(k(t)\) concurrently
delegated subproblems, the inference quality available to \(L_n\) for
genuinely novel problems satisfies:

\[\rho(L_n, a_{\text{novel}}) = \frac{|K^{[a_{\text{novel}}]}|}{|K^{[a_{\text{novel}}]}| + k(t)\cdot\alpha}\]

Top-down allocation is viable when \(k(t)\) is small relative to
\(C_n\). As the entity's informed action space \(A\) grows, \(k(t)\)
grows proportionally and \(\rho \to 0\) unless \(k(t)\) is bounded by
an independent mechanism. Escalation (Theorem 2) provides that
mechanism: by handling routine problems below \(L_n\), escalation keeps
\(k(t)\) bounded to genuinely novel concurrent subproblems, preserving
\(C_n\) capacity for the top-down allocation that those subproblems
require.

\emph{Proof.} \(L_n\) holds \(K^{[a_{\text{novel}}]}\) for the current
problem plus \(k(t)\cdot\alpha\) for delegation overhead. Since
\(\alpha > 0\) and \(k(t)\) grows with \(|A|\), the denominator is
unbounded while \(|K^{[a_{\text{novel}}]}|\) remains bounded (any
specific problem requires only its relevant subgraph). Therefore
\(\rho \to 0\) as \(|A| \to \infty\). When
\(k(t)\cdot\alpha \geq C_n - |K^{[a_{\text{novel}}]}|\), \(L_n\) lacks
capacity for genuine inference entirely. To maintain \(\rho\) at a fixed
level, \(C_n\) must grow as \(k(t)\) grows. \(\square\)

\emph{Note on attention mechanisms.} The \(O(k)\) delegation overhead
applies to any mechanism that evaluates all \(n\) stored propositions to
determine relevance, including full self-attention (\(O(n^2)\) per
query). Sparse and learned attention mechanisms reduce this cost by
learning fixed relevance patterns. See the remark on continuous-weight
mechanisms in Section 10 for a discussion of how these optimizations
relate to stored adjacency.

\emph{Note.} This result is consistent with the Free Energy Principle's
own minimization imperative {[}Friston2010, Friston2017{]}. FEP holds
that intelligent systems minimize the variance they must manage. Growing
\(C_n\) to accommodate unbounded delegation overhead adds variance: more
propositions in \(L_n\) means more uncertain relevance to any specific
problem. Escalation is the mechanism that keeps delegation overhead
bounded, preserving the conditions under which precision allocation
operates effectively. Without escalation, precision allocation is
self-undermining: allocating attention from \(L_n\) increases the
variance that \(L_n\) must hold in order to allocate.

\emph{Biological note.} The cortical hierarchy is bidirectional:
top-down signals exist at every level, and prefrontal cortex
demonstrably performs top-down allocation across cortical networks ---
biasing sensory processing, gating working memory, selecting task sets
{[}Friston2010{]}. This is not a counterexample to the theorem; it is
the theorem's prediction. Prefrontal top-down allocation operates at the
granularity of task-sets and functional networks, not individual
circuits. Its capacity is bounded at approximately four concurrent
allocations {[}Cowan2001{]} --- precisely the \(O(k)\) constraint the
theorem derives. Within each allocated network, processing is local and
automatic: escalation architecture at the circuit level. What triggers
prefrontal engagement is bottom-up conflict or novelty --- a lower
level's inability to bound uncertainty --- not a top-down scan of all
ongoing processes. Escalation determines engagement; top-down allocation
operates within the freed capacity. The transition from active to
automatic processing across skill acquisition is the biological
measurement: as \(f\) becomes automatic, it stops consuming prefrontal
capacity, keeping \(k(t)\) small enough that top-down allocation
remains viable for genuinely novel problems.

\textbf{Theorem 2b (Graph-Structured Knowledge Bounds Active Context).}
Under the escalation architecture (Theorem 2) with \(K\) structured as a
directed graph (Definition 2), as \(K\) grows:

\begin{enumerate}
\def\labelenumi{\arabic{enumi}.}
\tightlist
\item
  \(|\text{ctx}(L_n, t)| = |K^{[a]}|\) for each problem reaching
  \(L_n\): no delegation overhead
\item
  \(\rho(L_n, f) = 1\) for all problems at \(L_n\): inference quality is
  maximum
\item
  \(C_n\) need not grow as \(|A|\) or \(|K|\) grows
\end{enumerate}

\emph{Proof.}

\textbf{Part 1 (No delegation overhead).} By Theorem 2, \(L_n\) is
engaged only when all \(L_{i < n}\) fail. Problems handled below \(L_n\)
never appear in \(L_n\): they generate no delegation context there.
\(L_n\) receives a problem only when the escalation signal propagates up
from below --- not because \(L_n\) delegated and is monitoring, but
because the lower level failed and the failure is passed upward. The
distinction is causal: escalation is a push from failure, not a pull
from delegation. Therefore \(L_n\) holds no monitoring context for
subproblems in progress at lower levels.

\textbf{Part 2 (\(\rho = 1\) conditional on edge fidelity).} The graph's
edge structure encodes inferential and causal relevance between
propositions. Graph traversal from the escalated problem node follows
edges to exactly \(K^{[a]}\) --- the propositions inferentially
connected to \(f\) --- without loading any others.
\(|\text{ctx}(L_n, t)| = |K^{[a]}|\), so \(\rho = 1\) by Definition 8a.
\emph{This result is conditional on edge fidelity: \(\rho = 1\) holds
when the graph's edges correctly encode all and only the inferential
relevance of each proposition. Missing edges cause under-retrieval
(\(\rho\) remains 1 by count but \(K^{[a]}\) is incomplete); spurious
edges cause over-retrieval (\(\rho < 1\)). Graph structure in principle
supports \(\rho = 1\); achieving it requires well-specified edges.}

\textbf{Part 3 (Graph is the necessary structure under joint
constraints).} The escalation architecture and bounded context impose
three simultaneous constraints on any knowledge representation:

\begin{itemize}
\tightlist
\item
  \textbf{Retrieval}: the cost of accessing the knowledge relevant to
  problem \(f\) must be \(O(|K^{[a]}|)\) --- proportional to the
  subgraph needed, not to \(|K|\). As \(|K|\) grows, retrieval cost
  cannot grow with it, or \(L_n\) is burdened proportionally to past
  learning.
\item
  \textbf{Context}: \(|\text{ctx}(L_n, t)| = |K^{[a]}|\) exactly ---
  \(\rho = 1\). No extraneous propositions may enter \(L_n\); each one
  admitted that is not in \(K^{[a]}\) is noise that degrades inference
  (Theorem 2a).
\item
  \textbf{Memory}: shared propositions must be stored once. A
  proposition that grounds multiple informed actions across different
  \(f\) must exist as a single node, not duplicated at each point of
  use. Duplication means that when evidence updates a proposition, every
  copy must be updated --- an \(O(|A|)\) maintenance cost that grows
  without bound.
\end{itemize}

<<<<<<< HEAD
We claim any representation satisfying all three is a graph:

\emph{Proof by elimination.} A flat store violates constraint 1:
retrieving \(K^{[f]}\) requires scanning \(O(|K|)\) propositions to
identify the relevant ones. A tree violates constraint 3: a proposition
relevant to multiple \(f\) across different branches must appear in each
branch or the tree loses the cross-reference, duplicating nodes and
breaking shared-update consistency. A vector store (embedding-based
retrieval) violates constraint 2: approximate nearest-neighbor retrieval
includes false positives (propositions near \(K^{[f]}\) in embedding
space but not inferentially connected) and misses false negatives
(inferentially connected propositions that embed distantly), so
\(\rho < 1\) in general.

Any representation that satisfies all three must: (i) encode which
propositions connect to which in retrievable form --- so that traversal
of relevant propositions costs \(O(|K^{[f]}|)\), not \(O(|K|)\); and
(ii) store each proposition once with links to all contexts that use it
--- so that shared propositions remain consistent. A structure
satisfying both (i) and (ii) is, by definition, a directed graph with
typed edges. The graph is not the unique name for this structure, but it
is the unique structure satisfying the constraints. \(\square\)
=======
We claim graph structure is necessary: any representation that is not a
graph must lose information as new propositions are added.

\emph{Proof by information preservation.} \(C_n\) is bounded (Definition
8a). As new propositions arrive, \(|K|\) grows. When \(|K|\) reaches
\(C_n\), the representation is in the curation regime: adding any new
proposition \(p_{new}\) requires making space. There are exactly two
ways to make space:

\textbf{Option A: Discard an existing proposition.} Without factoring
--- without explicit relational structure --- the entity has no
mechanism to determine which existing propositions are redundant. Two
propositions \(p_1, p_2\) that share a latent factor appear as
independent entries; there is no stored evidence that they share
structure. The choice of which to discard is therefore made without
information about relational dependencies. Discarding \(p_{old}\) loses
the information it encodes, and potentially also loses the information
that \(p_{old}\) has any relation to other propositions --- a
second-order loss. This is information loss.

\textbf{Option B: Factor --- extract shared structure, reducing \(|K|\)
without discarding.} If \(p_1\) and \(p_2\) share a latent factor
\(f^*\), extracting \(f^*\) as a single named node with typed edges to
\(p_1\) and \(p_2\) reduces \(|K|\) by replacing redundant structure
with a single shared reference. The information in \(p_1\) and \(p_2\)
is preserved; the information about their shared grounding in \(f^*\) is
now explicitly stored as an edge. Space is created without loss.
Factoring also resolves the update problem: when new evidence revises
\(f^*\), the single node is updated once, and the revision propagates to
all propositions that reference it via edges --- \(O(1)\) instead of
\(O(|A|)\).

Factoring is graph construction. Naming \(f^*\), creating a node, and
storing typed edges from \(p_1\) and \(p_2\) to it is precisely the
operation that produces a directed graph.

\textbf{Option C: Lossy compression --- replace propositions with a
summary that occupies less space.} This is a genuine third option, but
it reduces to Option A or Option B. Compressing \(p_1\) and \(p_2\)
into a summary \(p_s\) that discards detail is a form of discarding:
the lost detail may ground other propositions or informed actions. If
\(p_1\) has dependents --- propositions or actions that reference it ---
then replacing \(p_1\) with \(p_s\) silently invalidates those
dependents unless the entity knows what depends on the lost detail. But
knowing what depends on \(p_1\) \emph{is} knowing \(p_1\)'s adjacency
--- its edges. Without stored adjacency, compression is blind
discarding: the entity cannot evaluate what breaks. With stored
adjacency, the entity can scope the damage, notify dependents, and
compress safely --- but storing adjacency is maintaining a graph.
Compression that preserves correctness requires the same relational
structure as factoring. Compression that ignores relational structure is
Option A with extra steps.

Any representation that makes space without factoring or correctly
scoped compression loses information; any representation that factors or
correctly scopes compression is constructing a graph.

\emph{Consequence for alternative representations.} A flat store cannot
factor: shared structure has no representation, so at \(C_n\) it must
discard. A tree loses cross-branch propositions when one branch is
pruned. A vector store holds shared factors implicitly in embedding
geometry but cannot factor without naming --- and naming is graph
construction. Every representation that preserves information as
\(|K|\) grows beyond \(C_n\) must factor. Every factoring operation
creates nodes with typed edges. The graph is what factoring is.
\(\square\)
>>>>>>> 319f447daa46f8c1e949fb7929535d1aa6396025

\textbf{Part 4 (\(C_n\) stays bounded).} By Corollary 6, proven informed
actions promote to lower encoding levels and leave \(L_n\). Each
promotion reduces future \(L_n\) load. As \(K\) grows, more problems are
handled below \(L_n\), not above. \(L_n\) stays occupied by genuinely
novel uncertainty --- the frontier moves outward, but the frontier's
size at any moment is bounded. \(C_n\) need not grow; \(K\) grows
without filling \(L_n\). \(\square\)

\textbf{Corollary 4a (Selection Pressure Produces Graph-Structured
Knowledge).} Two entities with equal \(|A|\) and \(|K|\): one with
top-down allocation (growing \(C_n\)), one with graph + escalation
(bounded \(C_n\)). The second has strictly higher \(\rho\) at lower
substrate cost. Intelligence per unit of active context --- the ratio
\(I(E)/C_n\) --- is higher for the graph-structured entity. By Theorem
2b Part 3, graph structure is not merely more efficient than
alternatives --- it is the necessary consequence of satisfying the
retrieval, context, and memory constraints jointly. Entities that grow
\(C_n\) to manage delegation overhead are selected against because they
violate the context constraint as \(|K|\) grows; entities using
approximate retrieval are selected against because they violate the
\(\rho = 1\) requirement; entities duplicating shared propositions are
selected against because maintenance cost grows with \(|A|\). Selection
pressure converges on the graph because the graph is the only structure
that does not eventually fail one of the three constraints.

\textbf{Corollary 4b (The Scaling Argument Against Context Window
Growth).} The prevailing approach to scaling intelligence --- growing
\(C_n\) directly --- is self-defeating on two grounds {[}Liu2024{]}.
First, by Theorem 2a, more context means more noise: for any specific
problem, \(\rho\) decreases as \(C_n\) fills with content not relevant
to that problem. More context does not produce better inference; it
produces higher variance that the same inference must manage. Second, by
Definition 7 (\(I(E) = \mathcal{A} \cdot
\eta_M\)), intelligence is the rate at which \(M\) improves \((A, K)\)
--- not the capacity of \(L_n\). A larger \(L_n\) does not increase
\(\eta_M\); it increases the noise environment in which \(M\) must
operate. The efficient path to growing intelligence is not to grow
\(C_n\) but to grow the encoded knowledge at lower levels --- filling
the graph, not the context window.

\textbf{Theorem 2c (Graph Necessity from Retention Dynamics).} Theorem
2b Part 3 establishes graph necessity from the joint constraints of
retrieval, context, and memory under bounded \(C_n\). This theorem
establishes an independent necessity: graph structure is the only
representation that supports sub-linear evaluation of the retention
criterion (Theorem 1) under continuous operation.
<<<<<<< HEAD

\emph{Statement.} Let \(K\) be a knowledge representation supporting the
retention criterion of Theorem 1. The retention criterion requires four
operations:

\begin{enumerate}
\def\labelenumi{(\roman{enumi})}
\tightlist
\item
  \textbf{Addition with provenance}: insert a proposition \(p\) together
  with its evidential grounds --- which existing propositions \(p\)
  depends on, and by what relation (inferential, causal, evidential)
\item
  \textbf{Contextual retrieval}: given an action context \(f\), retrieve
  the propositions relevant to \(f\)
\item
  \textbf{Dependency query}: given a proposition \(p\), determine which
  other propositions are grounded through \(p\) --- i.e., which
  propositions would require re-evaluation if \(p\) were removed or
  revised
\item
  \textbf{Selective removal with cascade}: when \(p\) fails the
  retention criterion (or is displaced in the curation regime), remove
  \(p\) and re-evaluate every proposition whose grounding depends on
  \(p\), transitively, until the knowledge base stabilizes
\end{enumerate}

For any knowledge representation that achieves sub-linear cost on all
four operations simultaneously, the representation necessarily encodes
directed adjacency between propositions --- and any structure encoding
directed adjacency over propositions with typed relations \emph{is} a
semantic graph, regardless of implementation.

\emph{Proof.}

\textbf{Part 1 (Dependency queries are necessary and continuous).} The
retention criterion of Theorem 1 operates continuously: every new piece
of evidence potentially shifts the cost-benefit of existing
propositions. When evidence revises or invalidates a proposition \(p\),
the entity must evaluate every proposition whose grounding passes
through \(p\). This is not a batch operation that can be deferred ---
acting on a proposition whose ground has been invalidated risks
\(U > U_{lethal}^d\) (Definition 5a). The dependency query ``what is
grounded through \(p\)?'' is therefore a survival-critical operation
invoked at the frequency of evidence arrival.

\textbf{Part 2 (Lower bound on non-adjacency dependency queries).}
Consider a representation \(R\) that stores \(n\) propositions without
explicit directed relationships between them (a flat store, hash table,
embedding space, or any structure that does not record which propositions
depend on which). To answer the dependency query for proposition \(p\),
\(R\) must examine each of the other \(n - 1\) propositions to determine
whether its grounding involves \(p\). Without stored adjacency, there is
no mechanism to restrict this examination: any proposition might depend
on \(p\), and the only way to determine this is to check.

The cost of a single dependency query is therefore \(\Omega(n)\) for any
representation without stored adjacency. Since dependency queries are
invoked at evidence-arrival frequency (Part 1), the amortized cost of
maintaining retention is \(\Omega(n)\) per evidence event --- growing
linearly with the size of \(K\).

\textbf{Part 3 (Cascade amplifies the cost).} Selective removal under
the retention criterion is not a single deletion. When \(p\) is removed,
each proposition \(q\) that was grounded through \(p\) must be
re-evaluated: does \(q\) have alternative grounds, or does it now fail
the retention criterion? If \(q\) also fails, its dependents must be
re-evaluated in turn. This is a transitive closure operation on the
dependency relation.

In a representation with explicit adjacency (a graph), cascade follows
edges: the cost is \(O(|D_p|)\) where \(|D_p|\) is the number of
propositions in the transitive dependency set of \(p\). In a
representation without adjacency, each level of the cascade requires an
\(\Omega(n)\) scan to identify dependents, and the cascade may propagate
through multiple levels. The total cost of a single removal with cascade
is \(\Omega(n \cdot L)\) where \(L\) is the cascade depth --- versus
\(O(|D_p|)\) in the graph, where \(|D_p| \ll n\) for typical
propositions with bounded connectivity.

\textbf{Part 4 (Any sub-linear solution encodes adjacency).} Suppose a
representation \(R\) achieves \(o(n)\) cost for dependency queries. Then
\(R\) must store, for each proposition, information about which other
propositions depend on it --- otherwise the \(\Omega(n)\) lower bound of
Part 2 applies. This stored information is a set of directed
relationships: for each proposition \(p\), a record of the propositions
that reference \(p\) as part of their grounding.

This is precisely a directed adjacency structure. The nodes are
propositions. The edges are directed dependency relations (inferential,
causal, evidential). The structure supports traversal from any node to
its dependents. Any representation that stores this information ---
regardless of whether it is implemented as adjacency lists, edge tables,
pointers, or any other mechanism --- \emph{is} a directed graph over
propositions.

\textbf{Part 5 (Semantic typing is necessary for correct cascade).} A
bare adjacency structure (propositions linked by untyped edges) is
insufficient for correct cascade evaluation. When \(p\) is invalidated,
the cascade must distinguish between propositions that are
\emph{inferentially} grounded through \(p\) (which require
re-evaluation) and propositions that are merely \emph{associated} with
\(p\) (which do not). An untyped edge conflates these: every neighbor of
\(p\) must be treated as potentially dependent, degrading the cascade to
the full neighborhood rather than the true dependency set.

Typed edges --- edges carrying the relation type (``grounds,''
``evidences,'' ``causes,'' ``supports'') --- enable the cascade to
follow only the dependency-bearing edges, keeping cascade cost
proportional to the true dependency set \(|D_p|\) rather than the full
neighborhood \(|\text{adj}(p)|\). Since \(|\text{adj}(p)|\) can be much
larger than \(|D_p|\) (a proposition may be associated with many others
without grounding them), typed edges are necessary for efficient cascade.

A directed graph with typed edges over propositions is a semantic graph.
\(\square\)

\emph{Consequence.} Theorem 2c establishes the complete set of
structural requirements for any knowledge representation supporting the
retention criterion of Theorem 1: (i) sub-linear contextual retrieval,
(ii) sub-linear dependency query, (iii) sub-linear selective removal
with cascade, and (iv) correct cascade scoping via typed relations.
These requirements are independent of Theorem 2b's static constraints
(retrieval, context, memory under bounded \(C_n\)) and provide a second,
orthogonal proof of graph necessity from the dynamics of continuous
knowledge maintenance.

\emph{Corollary 4c (Edge Growth Dominates Node Growth).} In a semantic
graph under the retention criterion, the ratio of edges to nodes
increases as \(K\) matures. Each new proposition added with provenance
creates at least one edge (to its evidential ground); propositions that
ground multiple other propositions accumulate incoming edges over time
without themselves growing in size. Edges are the cheapest structure in
the graph --- they are pointers, not content. The graph's information
density therefore improves as it matures: the relational structure
(edges) that enables sub-linear cascade grows faster than the content
(nodes) that the structure organizes. This is the operational advantage
of the graph: the structure needed for maintenance is the cheapest
element to store, and it grows precisely where it is most needed ---
at high-dependency nodes whose removal would trigger the largest
cascades.

\begin{center}\rule{0.5\linewidth}{0.5pt}\end{center}
=======
>>>>>>> 319f447daa46f8c1e949fb7929535d1aa6396025

\emph{Statement.} Let \(K\) be a knowledge representation supporting the
retention criterion of Theorem 1. The retention criterion requires four
operations:

<<<<<<< HEAD
\textbf{Theorem 3 (Agency Necessarily Produces Gradient Descent).}

\emph{On the space of descent.} The gradient descent in this theorem
operates in \(F\) --- the space of informed actions --- not in \(W\).
This distinction is critical. \(W\) may be arbitrarily rough,
discontinuous, or hostile; no smoothness assumption on \(W\) is required
or implied. The entity does not descend over \(W\); it descends over its
own function space, updating \(f \in F\) in response to signals received
from \(W\).

\(F \subseteq (K \times W_{obs} \to \Delta A)\) maps into \(\Delta A\),
the space of probability distributions over actions. Distribution spaces
are convex and metrizable (under KL divergence, total variation, or
Wasserstein distance); they support well-behaved subgradient structure
regardless of the topology of the domain \(W\) from which signals
arrive. The smoothness conditions for convergence are conditions on
\(F\), not on \(W\): specifically, that \(U\) as a functional on \(F\)
is locally Lipschitz, which holds whenever the likelihood factors
\(\ell_p\) (Definition 5) are bounded and measurable --- a minimal
requirement for any well-specified knowledge graph. Under these
conditions on \(F\), stochastic subgradient descent converges in the
function space independently of whether \(W\) is smooth.

\emph{Continuity note.} Since \(F\) is defined over \(W\), no
\(f \in F\) can produce a world state outside \(W\). Combined with the
continuity of time, the trajectory \(\tau(t)\) through \(W \times T\) is
continuous, \(U(\tau(t))\) is continuous in \(t\), and \(\frac{dU}{dt}\)
exists almost everywhere. The gradient \(\nabla_F U\) is this temporal
derivative, well-defined by the smoothness of \(F\) as a function space,
not by any assumption on \(W\).

\emph{Part I: Compulsion.} When an informed action \(f_i\) fails in
world state \(w\), \(W\) returns a feedback signal
\(\delta = \phi(f_i, w) - \mathbb{E}[\phi(f_i, \cdot) \mid K]\)
(Definition 3a): the signed deviation of observed outcome from
expectation. This signal arrives from \(W\) but is applied in \(F\): it
is a subgradient of \(U\) as a functional on \(F\) at the current
\(f_i\), evaluated at the encountered \(w\). The encountered world
states are on-policy --- determined by the entity's trajectory, not
sampled i.i.d. from \(P(w \mid K)\) --- but this does not impede
convergence. The subgradient \(\delta\) is valid at every encountered
\(w\) regardless of how \(w\) was reached; the entity is updating a
function in \(F\), and each \(\delta\) is a legitimate descent direction
in that space. The burn tells you that this \(f_i\) was insufficient at
\(w\). The entity does not choose whether to receive this signal. \(W\)
imposes it.
=======
\begin{enumerate}
\def\labelenumi{(\roman{enumi})}
\tightlist
\item
  \textbf{Addition with provenance}: insert a proposition \(p\) together
  with its evidential grounds --- which existing propositions \(p\)
  depends on, and by what relation (inferential, causal, evidential)
\item
  \textbf{Contextual retrieval}: given an action context \(f\), retrieve
  the propositions relevant to \(f\)
\item
  \textbf{Dependency query}: given a proposition \(p\), determine which
  other propositions are grounded through \(p\) --- i.e., which
  propositions would require re-evaluation if \(p\) were removed or
  revised
\item
  \textbf{Selective removal with cascade}: when \(p\) fails the
  retention criterion (or is displaced in the curation regime), remove
  \(p\) and re-evaluate every proposition whose grounding depends on
  \(p\), transitively, until the knowledge base stabilizes
\end{enumerate}

For any knowledge representation that achieves sub-linear cost on all
four operations simultaneously, the representation necessarily encodes
directed adjacency between propositions --- and any structure encoding
directed adjacency over propositions with typed relations \emph{is} a
semantic graph, regardless of implementation.

\emph{Proof.}

\textbf{Part 1 (Dependency queries are necessary and continuous).} The
retention criterion of Theorem 1 operates continuously: every new piece
of evidence potentially shifts the cost-benefit of existing
propositions. When evidence revises or invalidates a proposition \(p\),
the entity must evaluate every proposition whose grounding passes
through \(p\). This is not a batch operation that can be deferred ---
acting on a proposition whose ground has been invalidated risks
\(U > U_{lethal}^d\) (Definition 5a). The dependency query ``what is
grounded through \(p\)?'' is therefore a survival-critical operation
invoked at the frequency of evidence arrival.

\textbf{Part 2 (Lower bound on non-adjacency dependency queries).}
Consider a representation \(R\) that stores \(n\) propositions without
explicit directed relationships between them (a flat store, hash table,
embedding space, or any structure that does not record which propositions
depend on which). To answer the dependency query for proposition \(p\),
\(R\) must examine each of the other \(n - 1\) propositions to determine
whether its grounding involves \(p\). Without stored adjacency, there is
no mechanism to restrict this examination: any proposition might depend
on \(p\), and the only way to determine this is to check.

The cost of a single dependency query is therefore \(\Omega(n)\) for any
representation without stored adjacency. Since dependency queries are
invoked at evidence-arrival frequency (Part 1), the amortized cost of
maintaining retention is \(\Omega(n)\) per evidence event --- growing
linearly with the size of \(K\).

\textbf{Part 3 (Cascade amplifies the cost).} Selective removal under
the retention criterion is not a single deletion. When \(p\) is removed,
each proposition \(q\) that was grounded through \(p\) must be
re-evaluated: does \(q\) have alternative grounds, or does it now fail
the retention criterion? If \(q\) also fails, its dependents must be
re-evaluated in turn. This is a transitive closure operation on the
dependency relation.

In a representation with explicit adjacency (a graph), cascade follows
edges: the cost is \(O(|D_p|)\) where \(|D_p|\) is the number of
propositions in the transitive dependency set of \(p\). In a
representation without adjacency, each level of the cascade requires an
\(\Omega(n)\) scan to identify dependents, and the cascade may propagate
through multiple levels. The total cost of a single removal with cascade
is \(\Omega(n \cdot L)\) where \(L\) is the cascade depth --- versus
\(O(|D_p|)\) in the graph, where \(|D_p| \ll n\) for typical
propositions with bounded connectivity.

\textbf{Part 4 (Any sub-linear solution encodes adjacency).} Suppose a
representation \(R\) achieves \(o(n)\) cost for dependency queries. Then
\(R\) must store, for each proposition, information about which other
propositions depend on it --- otherwise the \(\Omega(n)\) lower bound of
Part 2 applies. This stored information is a set of directed
relationships: for each proposition \(p\), a record of the propositions
that reference \(p\) as part of their grounding.

This is precisely a directed adjacency structure. The nodes are
propositions. The edges are directed dependency relations (inferential,
causal, evidential). The structure supports traversal from any node to
its dependents. Any representation that stores this information ---
regardless of whether it is implemented as adjacency lists, edge tables,
pointers, or any other mechanism --- \emph{is} a directed graph over
propositions.

\textbf{Part 5 (Semantic typing is necessary for correct cascade).} A
bare adjacency structure (propositions linked by untyped edges) is
insufficient for correct cascade evaluation. When \(p\) is invalidated,
the cascade must distinguish between propositions that are
\emph{inferentially} grounded through \(p\) (which require
re-evaluation) and propositions that are merely \emph{associated} with
\(p\) (which do not). An untyped edge conflates these: every neighbor of
\(p\) must be treated as potentially dependent, degrading the cascade to
the full neighborhood rather than the true dependency set.

Typed edges --- edges carrying the relation type (``grounds,''
``evidences,'' ``causes,'' ``supports'') --- enable the cascade to
follow only the dependency-bearing edges, keeping cascade cost
proportional to the true dependency set \(|D_p|\) rather than the full
neighborhood \(|\text{adj}(p)|\). Since \(|\text{adj}(p)|\) can be much
larger than \(|D_p|\) (a proposition may be associated with many others
without grounding them), typed edges are necessary for efficient cascade.

A directed graph with typed edges over propositions is a semantic graph.
\(\square\)

\emph{Consequence.} Theorem 2c establishes the complete set of
structural requirements for any knowledge representation supporting the
retention criterion of Theorem 1: (i) sub-linear contextual retrieval,
(ii) sub-linear dependency query, (iii) sub-linear selective removal
with cascade, and (iv) correct cascade scoping via typed relations.
These requirements are independent of Theorem 2b's static constraints
(retrieval, context, memory under bounded \(C_n\)) and provide a second,
orthogonal proof of graph necessity from the dynamics of continuous
knowledge maintenance.

\emph{Corollary 4c (Edge Growth Dominates Node Growth).} In a semantic
graph under the retention criterion, the ratio of edges to nodes
increases as \(K\) matures. Each new proposition added with provenance
creates at least one edge (to its evidential ground); propositions that
ground multiple other propositions accumulate incoming edges over time
without themselves growing in size. Edges are the cheapest structure in
the graph --- they are pointers, not content. The graph's information
density therefore improves as it matures: the relational structure
(edges) that enables sub-linear cascade grows faster than the content
(nodes) that the structure organizes. This is the operational advantage
of the graph: the structure needed for maintenance is the cheapest
element to store, and it grows precisely where it is most needed ---
at high-dependency nodes whose removal would trigger the largest
cascades.

\begin{center}\rule{0.5\linewidth}{0.5pt}\end{center}

\subsection{6. Gradient Descent on
Uncertainty}\label{gradient-descent-on-uncertainty}

\textbf{Theorem 3 (Agency Necessarily Produces Evaluation-Driven State
Revision).}

\emph{On the space of descent and dimensionality.} The gradient descent
in this theorem operates in \(A\) --- the space of informed actions ---
not in \(W\). \(W\) may be arbitrarily rough, discontinuous, or of
unknown dimensionality; no assumption on \(W\)'s structure is required.
The entity does not descend over \(W\); it updates \(a \in A\) in
response to signals received from \(W\).

The dimensionality of \(W\) or of \(M_{\text{res}}\) is not assumed. We
do not know whether the latent space is finite or infinite-dimensional.
The convergence proof does not require this knowledge. It operates on
the scalar \(U(w, K) = H(w \mid K)\), which is bounded below by zero
regardless of the dimensionality of \(W\) or \(A\). The proof shows
\(\mathbb{E}[U]\) decreases monotonically under \(M\)-updates and
converges by the monotone convergence theorem, without decomposing \(U\)
into independent dimensions.

\emph{Part I: Forced Signal, Response Gate.} Every executed action
\(a_i\) in world state \(w\) generates a feedback signal
\(\delta = \phi(a_i, w) - \mathbb{E}[\phi(a_i, \cdot) \mid K]\)
(Definition 3a): the signed deviation of observed outcome from
expectation. This signal arrives from \(W\) but is applied in \(A\): it
is a subgradient of \(U\) as a functional on \(A\) at the current
\(a_i\), evaluated at the encountered \(w\). The signal is generated by
execution, not by \(\mathcal{A}\). It arrives whether or not the entity
has \(\mathcal{A}\). The entity does not choose whether to receive this
signal. \(W\) imposes it.

\(\mathcal{A}\) is the response gate: the probability that \(M\) engages
on \(\delta\) when \(\delta\) arrives. An entity with
\(\mathcal{A} = 0\) receives every signal but \(M\) does not engage;
\((A, K)\) remains unchanged. An entity with \(\mathcal{A} > 0\) engages
\(M\) on the signal, updating \((A, K)\) in response. The encountered
world states are on-policy --- determined by the entity's trajectory,
not sampled i.i.d. from \(P(w \mid K)\) --- but this does not impede
convergence. The subgradient \(\delta\) is valid at every encountered
\(w\) regardless of how \(w\) was reached; the entity is updating a
function in \(A\), and each \(\delta\) is a legitimate descent direction
in that space. The burn tells you that this \(a_i\) was insufficient at
\(w\).
>>>>>>> 319f447daa46f8c1e949fb7929535d1aa6396025

An entity with \(\mathcal{A}\) that does not update in response to
\(\delta\) does not reduce \(U\). If \(U\) does not decrease, it
eventually exceeds \(U_{lethal}^d\) in some critical domain. The entity
does not survive. Selection removes it. Therefore any entity with
\(\mathcal{A}\) that persists must follow \(\delta\):

<<<<<<< HEAD
\[\Delta f_i = -\eta \cdot \delta\]

This is stochastic subgradient descent in \(F\): each step follows a
valid subgradient of \(U\) as a functional on the function space,
evaluated at an encountered world state. The on-policy nature of the
trajectory does not introduce bias into the subgradient direction; it
affects the distribution of encountered states but not the validity of
each \(\delta\) as a descent direction. Under the Lipschitz condition on
\(F\) stated above, the sequence \(\{f_i\}\) converges to a local
minimum of \(U\) in \(F\). This is not chosen. Any entity with
\(\mathcal{A} > 0\) that persists must have followed \(\delta\): the
gradient is not a strategy but what survival-filtered updating is.
=======
\[\Delta a_i = -\eta \cdot \delta\]

\emph{Convergence of U.} The convergence argument does not decompose
by dimension. \(U(w, K) = H(w \mid K) \geq 0\) (Definition 5) is a
scalar quantity bounded below by zero. We show it decreases in
expectation under \(M\)-updates and use this to establish convergence
directly.

\textbf{Step 1 (Each update reduces U in expectation).} By Definition
3b, every informed action \(a \in A\) applied at world state \(w\) yields
information gain \(G(a, K) = I(a\,;\,w \mid K) \geq 0\), with equality
only when \(a\) provides no new information about \(w\) given \(K\).
When \(M\) engages on the feedback signal \(\delta\) and updates
\((A, K) \to (A', K')\), the updated knowledge \(K'\) satisfies
\(H(w \mid K') \leq H(w \mid K)\) --- conditioning on additional
evidence cannot increase entropy (data processing inequality,
{[}Cover2006{]}). Therefore:

\[\mathbb{E}[U(w, K') \mid K] \leq U(w, K) - G(f, K)\]

Each \(M\)-update that engages on a signal with \(G > 0\) strictly
reduces expected uncertainty. Updates with \(G = 0\) leave \(U\)
unchanged.

\textbf{Step 2 (Convergence).} The sequence
\(\{\mathbb{E}[U(w, K(t))]\}_{t \geq 0}\) is non-increasing (by Step 1)
and bounded below by zero. By the monotone convergence theorem,
\(\mathbb{E}[U]\) converges. Since
\(\sum_t G(a_t, K(t)) \leq U(w, K(0)) < \infty\), the information gains
\(G(a_t, K(t)) \to 0\): the marginal improvement per update vanishes
asymptotically. This is convergence to a local minimum of \(U\) given
the entity's current \(A\) and \(\sigma\).

\textbf{Step 3 (Dimensionality is irrelevant to convergence).} The
argument above operates on the scalar \(U\), not on individual
dimensions of \(W\) or components of \(A\). The entity's
parameterization of \(A\) may grow as
\(\mathcal{T}_{reachable}\) expands and new propositions are grounded
(Part II). This growth does not invalidate convergence: each new
proposition \(p_{new}\) added to \(K\) can only reduce \(U\) (by the
data processing inequality), and the cross-grounding of \(p_{new}\) with
existing propositions can reduce \(U\) further by improving the
resolution of knowledge already in \(K\). Adding structure to \(K\) is
analogous to adding a column to a matrix and re-normalizing: the
existing entries can become better resolved, not worse. The coupling
between new and existing knowledge \emph{accelerates} convergence; it
cannot prevent it, because the Lyapunov bound \(U \geq 0\) holds
regardless of dimensionality.
>>>>>>> 319f447daa46f8c1e949fb7929535d1aa6396025

\emph{Part II: Non-termination.} Let \(\mathcal{T}_{reachable}(K)\) be
the set of trajectories reachable from the entity's current state given
knowledge \(K\) (as defined in Theorem 1). We claim that for any finite
\(K(t)\), the frontier of \(\mathcal{T}_{reachable}\) contains world
states \(w\) where \(H(w \mid K) > 0\), provided \(W\) is sufficiently
rich that each executed action can ground a genuinely new proposition.
In a finite \(W\) where all propositions are exhaustible, the frontier
could in principle be covered; the claim applies in any \(W\) where the
entity's epistemic horizon has not been closed --- which is the intended
domain.

By Definition 3a, executing any \(a \in A\) produces feedback
\(\phi(a, w)\) that grounds a new proposition \(p_a\) not previously in
\(K\), yielding \(K_a = K \cup \{p_a\}\). This richer \(K_a\) grounds
new informed actions not available from \(K\) alone (Definitions 3 and
<<<<<<< HEAD
3b): actions that require \(p_f\) to select or apply, reaching world
states not reachable without \(p_f\). Therefore
\(\mathcal{T}_{reachable}(K_f) \supseteq
\mathcal{T}_{reachable}(K)\), strictly whenever \(p_f\) grounds at least
one new \(f' \in F\). Since Part I establishes that gradient descent
produces feedback at every step, \(\mathcal{T}_{reachable}\) grows
=======
3b): actions that require \(p_a\) to select or apply, reaching world
states not reachable without \(p_a\). Therefore
\(\mathcal{T}_{reachable}(K_a) \supseteq
\mathcal{T}_{reachable}(K)\), strictly whenever \(p_a\) grounds at least
one new \(a' \in A\). Since Part I establishes that every executed
action generates a feedback signal, \(\mathcal{T}_{reachable}\) grows
>>>>>>> 319f447daa46f8c1e949fb7929535d1aa6396025
monotonically with \(K\). At any finite time \(t\), \(K(t)\) has not
covered \(\mathcal{T}_{reachable}(K(t))\): the frontier always contains
states \(w\) where \(H(w \mid K) > 0\). Therefore \(\nabla_A U\) is
never globally zero over \(\mathcal{T}_{reachable}\), and descent never
terminates within the entity's reachable epistemic domain.

<<<<<<< HEAD
\emph{Proof.} Part I: follows from \(U_{lethal}^d\) and selection;
entities that do not update are removed. Part II: follows from the
monotonic growth of \(\mathcal{T}_{reachable}\) with \(K\) --- each
executed action extends the reachable frontier, ensuring new unexplored
states are always accessible within the domain of the entity's possible
experience. \(\square\)

\textbf{Corollary 5 (Gradient Descent is Survival, Not Metaphor).} For
individual learning, gradient descent on \(U\) is not a description of
how learning happens to work. It is what survival-driven updating
\emph{is}: the entity that burns and does not become more careful does
not survive. The entity that survives is, by definition, the one that
followed the gradient. At the evolutionary timescale, the same
directional imperative operates on the population distribution over
\(F\)-variants --- not as an identical process but as a formally
analogous one, filtering for variants that reduce \(\mathbb{E}[U]\)
across the distribution of environments the lineage encounters.

\textbf{Remark (Multi-timescale Gradient Structure).} The formal
structure of individual learning is analogous to population-level
selection in the following sense: both can be described as descent on an
uncertainty objective, differing in the averaging domain and the object
that changes. This is a structural analogy, not a theorem --- the two
processes operate on mathematically distinct objects and the analogy
does not constitute an identity. It is recorded here because it
illuminates the relationship between learning rates at different
timescales, from which a provable consequence follows (see Derivation,
Part II).

\[\Delta A_{\text{individual}} = -\eta_{\text{ind}} \nabla_A \, U(w, K)\]

\[\Delta F_{\text{species}} \sim -\eta_{\text{evo}} \nabla_F \, \mathbb{E}_{\text{individuals}}\left[U(w, K)\right]\]

where \(\eta_{\text{evo}} \ll \eta_{\text{ind}}\) and \(\sim\) denotes
formal analogy rather than identity.

\emph{Derivation.}

\textbf{Part I (Analogous structure, different mathematical types).}
Individual learning updates \(F\) for a single entity based on
\(U(w, K)\) at specific world states \(w\) encountered by that entity.
Evolutionary change operates on a different object: it updates the
population-level distribution over \(F\)-variants through differential
reproduction, shifting mean fitness across a population of entities.
These are not the same mathematical operation --- individual learning
changes a function within one agent; evolutionary change changes a
probability distribution over many agents. The analogy is that both can
be described as moving in the direction that reduces the relevant
uncertainty objective: \(U(w, K)\) for the individual,
\(\mathbb{E}[U(w, K)]\) averaged across individuals for the population.
The mechanisms are distinct --- synaptic change vs.~differential
reproduction --- and the objects are distinct --- one entity's \(F\)
vs.~a population distribution over \(F\)-variants --- but the
directional relationship to \(U\) is formally analogous. The notation
\(\Delta F_{\text{species}}\) should be read as the population-level
analog of \(\Delta F\), not as the same operation on the same object.
=======
\emph{Proof.} Part I: \(U \geq 0\) is a Lyapunov function; each
\(M\)-update with \(G > 0\) strictly reduces \(\mathbb{E}[U]\) (data
processing inequality); \(\mathbb{E}[U]\) converges by monotone
convergence. Entities with \(\mathcal{A} = 0\) do not engage \(M\) and
are removed by selection. Part II: follows from the monotonic growth of
\(\mathcal{T}_{reachable}\) with \(K\) --- each executed action extends
the reachable frontier, ensuring new unexplored states are always
accessible within the domain of the entity's possible experience.
\(\square\)

\textbf{Corollary 5 (Evaluation-Driven State Revision is Survival, Not
Metaphor).} For individual learning, gradient descent on \(U\) is not a
description of how learning happens to work. It is what survival-driven
updating \emph{is}: the entity that burns and does not become more
careful does not survive. The entity that survives is, by definition,
the one whose \(M\) engaged on the signal. At the evolutionary
timescale, the mechanism differs in kind: random mutation perturbs the
\(L_0\) encoding, generating new \((A, K, M, \mathcal{A})\)
configurations; \(W\) evaluates each through the survival criterion;
configurations that persist reproduce. No individual organism follows a
gradient --- the gradient is revealed by differential survival of
randomly generated candidates. This is structurally the same as a
locally irrational action (temporarily negative \(\eta_M\)) at the
individual level: a perturbation that generates a new starting point for
the \(M \times \mathcal{A} \times K\) cross-product to be evaluated by
\(W\). The two mechanisms occupy a spectrum: at one end, random mutation
with binary survival feedback (evolution); at the other, directed update
via continuous \(\delta\) carrying full gradient information (individual
learning). Gradient descent is the efficient limit of this spectrum
under two conditions: (i) \(\delta\) carries directional information in
\(A\) --- not merely pass/fail, but a signed deviation identifying which
aspect of \(a_i\) was insufficient at \(w\); and (ii) \(M\) can use that
directional information to bias the update toward lower \(U\) rather
than generating a random candidate. When both conditions hold, each
evaluation step directly reduces \(U\). When neither holds, the
mechanism is random search with \(W\)-evaluation. Real learning systems
occupy positions between these limits; signal informativeness determines
how closely the process approximates directed gradient descent.

\textbf{Remark (Multi-timescale Structure of Evaluation-Driven
Revision).} Individual learning and evolutionary selection share a
common structure: execute, receive signal, retain or revise. The
mechanisms and mathematical objects differ in kind.

For individual learning, the mechanism is directed: the gradient signal
\(\delta\) arrives from a specific failure, \(M\) updates \((A, K)\) in
the direction that reduces \(U\):

\[\Delta A_{\text{individual}} = -\eta_{\text{ind}} \nabla_A \, U(w, K)\]

For evolutionary selection, the mechanism is undirected: random mutation
perturbs the \(L_0\) encoding of \((A, K, M, \mathcal{A})\), generating
candidate configurations. \(W\) evaluates each through survival. The
aggregate population-level distribution shifts toward configurations
with lower \(\mathbb{E}[U]\), but no individual organism descends a
gradient --- the gradient is revealed by differential survival of
randomly generated candidates. The notation
\(\Delta A_{\text{species}}\) denotes the shift in the population
distribution over \(A\)-variants, not a gradient step by any individual:

\[\Delta A_{\text{species}} = \operatorname{select}(\operatorname{mutate}(A_{\text{population}}))\]

Both processes are instances of the same structure: generate a candidate
state, submit to \(W\)'s evaluation, retain if valid. Gradient descent
(\(\Delta A_{\text{individual}}\)) is the efficient limit when signals
are informative enough to direct updates. Mutation followed by selection
(\(\Delta A_{\text{species}}\)) is the same mechanism without the
directed signal. This structural equivalence is not an identity ---
mechanisms, timescales, and mathematical objects are all distinct ---
but it grounds the common observation that both processes move toward
lower \(U\) at their respective timescales, and grounds a provable
consequence in Part II.

\emph{Derivation.}

\textbf{Part I (Common structure, different mechanisms).} Individual
learning updates \(A\) for a single entity based on \(U(w, K)\) at
specific world states \(w\) encountered by that entity. Evolutionary
change operates on a different object: it updates the population-level
distribution over \(A\)-variants through differential reproduction of
randomly generated candidates. These are not the same mathematical
operation --- individual learning changes a function within one agent
via directed update; evolutionary change changes a probability
distribution over many agents via undirected mutation and selection. The
common structure is: generate a new internal state, submit it to \(W\)'s
evaluation, retain if valid. Gradient descent is the efficient limit of
this structure; random mutation followed by selection is the general
case.
>>>>>>> 319f447daa46f8c1e949fb7929535d1aa6396025

\textbf{Part II (\(\eta_{\text{evo}} \ll \eta_{\text{ind}}\) follows
from Theorem 4).} The learning rate \(\eta_i\) at encoding level \(L_i\)
is inversely proportional to the rigor threshold \(\theta_i\): a higher
threshold requires more evidence per unit update, which is the
definition of a lower learning rate. Formally:

\[\eta_i \propto \frac{1}{\theta_i} = \frac{1}{1 - \varepsilon/C_i} = \frac{C_i}{C_i - \varepsilon}\]

For individual learning at \(L_n\): \(C_n\) is the cost of a single
wrong action, relatively small, so \(\varepsilon/C_n\) is
non-negligible, \(\theta_n = 1 - \varepsilon/C_n\) is meaningfully less
than 1, and \(\eta_{\text{ind}} = C_n/(C_n - \varepsilon)\) is
correspondingly large. For evolutionary learning at \(L_0\): \(C_0\) is
lineage survival, far exceeding \(\varepsilon\), so
\(\theta_0 = 1 - \varepsilon/C_0 \approx 1\) and
\(\eta_{\text{evo}} \approx 1\). Since
\(\theta_0 \approx 1 \gg \theta_n\), it follows that
\(1/\theta_0 \ll 1/\theta_n\), therefore
\(\eta_{\text{evo}} \ll \eta_{\text{ind}}\). The learning rate
difference is not assumed; it is derived from the cost-of-failure
difference established in Theorem 4.

An individual may establish an informed action that reduces \(U\) in one
context but fails in others. Only informed actions that reduce \(U\)
across sufficient variation in \(W\) survive the evolutionary test ---
the higher rigor threshold filters for robustness. The analogy captures
this shared directional property: both processes favor what reduces
\(U\) at the relevant scale. The timescale, the rigor threshold, the
object being updated, and the mechanism all differ; the directional
relationship to \(U\) is what the analogy preserves.

\textbf{Theorem 4 (Encoding Permanence Scales with Rigor).} The
certainty threshold required to promote an informed action \(f\) to
encoding level \(L_i\) is:

\[\theta_i = 1 - \frac{\varepsilon_{\text{acceptable}}}{C_i}\]

where \(C_i\) is the cost of failure at level \(L_i\) and
\(\varepsilon_{\text{acceptable}}\) is the maximum acceptable expected
cost of an encoding error, expressed in the same units as \(C_i\). For
heritable encoding (\(L_0\)), \(C_0\) is lineage survival, requiring
near-certainty across the full distribution of environments. For active
reasoning (\(L_n\)), \(C_n\) is the cost of a single wrong action, a far
lower bar.

\emph{Proof.} Encoding \(f\) at level \(L_i\) is irreversible at cost
\(r_i > 0\) (Definition 8). The expected cost of encoding an action with
current confidence \(\theta\) is:

\[\mathbb{E}[\text{encoding error cost}] = (1 - \theta) \cdot C_i\]

where \((1-\theta)\) is the probability the action is wrong and \(C_i\)
is the cost when it is. Promotion is safe when this expected cost does
not exceed what is acceptable:

\[(1 - \theta) \cdot C_i \leq \varepsilon_{\text{acceptable}}\]

Rearranging:

\[\theta \geq 1 - \frac{\varepsilon_{\text{acceptable}}}{C_i}\]

The minimum threshold at which promotion is safe is therefore
\(\theta_i = 1 - \varepsilon_{\text{acceptable}}/C_i\). Below this
threshold, the expected cost of encoding error exceeds
\(\varepsilon_{\text{acceptable}}\) and continued active engagement is
preferred. At or above it, encoding is the lower-expected-cost choice.
\(\square\)

\emph{Consequence.} You must be very certain to permanently change the
source code. A bad encoding at \(L_0\) does not kill one individual; it
kills the lineage. The rigor of the test is proportional to the
permanence of the consequence.

\emph{Note on error propagation through the hierarchy.} A potential
objection is that the proof derives \(\theta_i\) for each level
independently, whereas errors in a hierarchy compound: a failure at
\(L_0\) propagates upward. This objection assumes a naive cascading
model in which errors at \(L_i\) arrive at \(L_{i+1}\) at rate
\((1-\theta_i)\) regardless of whether \(L_i\) could address them. That
is not the escalation architecture.

Under escalation (Theorem 2), errors partition into two classes at each
level: \textbf{handled errors} --- failures \(L_i\) detects and resolves
within its scope, which never reach \(L_{i+1}\) --- and
\textbf{unhandled exceptions} --- failures outside \(L_i\)'s domain or
beyond its capacity to resolve, which escalate. The rate seen by
\(L_{i+1}\) is not the compounded Bernoulli rate from below. It is the
rate of genuinely irreducible errors at \(L_i\): errors that \(L_i\) was
correct not to handle, because the knowledge or capacity to resolve them
does not exist at \(L_i\). This is a different quantity from
\((1-\theta_i)\), and it is bounded by the design of \(\theta_i\) itself
--- the threshold ensures \(L_i\) handles everything within its
cost-justified scope.

Furthermore, when an unhandled exception escalates, the context at
\(L_i\) was already disrupted by the failure: the action failed, and
\(L_i\)'s current trajectory is suspended or invalid. The attention cost
at \(L_{i+1}\) is not paid on top of productive work; it is paid instead
of broken work. Escalation is correct triage. The thresholds
\(\theta_i\) are therefore independent by construction: each level's
\(\varepsilon_{\text{acceptable}}\) is its own irreducible error domain,
not an accumulation from below. The proof applies at each level without
modification.

\textbf{Definition 10 (Certain).} An informed action \(f\) is
\textbf{certain} when \(U(w, K_f) \approx 0\) across sufficient
variation in \(W\) to meet the promotion threshold \(\theta_i\) (Theorem
4). Certainty is asymptotic: decreasing density of uncertainty until
active engagement no longer justifies the residual gradient.

\textbf{Corollary 6 (Bounded Adaptation).} Certain informed actions
escape the active context window permanently, freeing \(L_n\) for the
next frontier of uncertainty.

\begin{center}\rule{0.5\linewidth}{0.5pt}\end{center}

\subsection{7. The Higher-Order Function
Space}\label{the-higher-order-function-space}

\textbf{Definition 11 (Higher-Order Function Space).} Let \(M\) be a
function space operating on the joint \((A, K)\) pair {[}Thrun1998{]}:

\[M : (A, K) \times \mathcal{L} \rightarrow (A, K)\]

where \(\mathcal{L}\) is a learning signal (experience, execution
results, selection pressure). \(M\) is the mechanism by which
\((A, K)\) changes --- categorically different from \(A\), not a more
certain version of it.

\(M\)'s type signature is complete without enumerating its internal
structure, because \(M\) is the cross product of all dimensions that
have not yet precipitated out at the current scale of the entity:

\[M = \prod_{i \in \text{remaining}} D_i\]

where the product ranges over all dimensions --- drive, meta-learning
rate, exploration policy, convergence criteria, and whatever else has
not yet reached the scale threshold to become a named dimension. \(A\),
\(K\), and \(\mathcal{A}\) are factors that have already been extracted
from this product. \(M\) is always the remainder: the undifferentiated
joint space of everything not yet named. As the entity scales,
additional dimensions precipitate out of \(M\) by the same normalization
argument that precipitated \(K\) from \(A\): when the shared structure
of a dimension across \(M\)-instances is large enough that factoring it
out reduces retrieval cost from \(O(n \cdot k)\) to \(O(n + k)\),
selection pressure forces the extraction. Each precipitation shrinks
\(M\) by one factor and names a new dimension.

The hierarchy is open upward: there is always potentially a \(M^{(2)}\)
that operates on \(M^{(1)}\) the same way \(M^{(1)}\) operates on
\((A, K)\). One level suffices until it does not --- when the gradient
can no longer be followed by improving \((A, K)\) alone, pressure
propagates upward. The required number of levels is determined by the
problem, not prescribed by the framework. The precipitation argument
applies at every level; the hierarchy is not a prescribed structure but
the output of a single recursive process applied at increasing scale.
See Open Question 7.

The degenerate cases of level merger generalize from the \(A = K\) case:
for any contiguous block of \(N\) precipitated dimensions together with
the \(M\) residual, all \(N+1\) collapse to a single undifferentiated
process when the domain is fully crystallized at all scales in that
block. The merger condition is \(U(w, K) \approx 0\) across the domain
at all \(N\) scales simultaneously --- one condition, not \(N\) pairwise
conditions. Since \(M\) represents the full cross product of remaining
dimensions, any merger including \(M\) collapses not just \(N\) named
levels but the entire unprecipitated hierarchy simultaneously. \(A = M\)
is the maximally primitive state: action spans the full cross product of
all remaining dimensions. \(M = M^{(2)}\) means the cross product is
self-similar under the meta-operation --- evolution approaches this
limit at the timescale of selection on selection mechanisms.

\textbf{Definition 11a (M Efficiency).} The efficiency \(\eta_M\) of
\(M\) is the rate at which \(M\) improves the joint \((A, K)\)
capability per unit of learning signal:

\[\eta_M(E) = \mathbb{E}_{\mathcal{L}} \left[ \frac{\max_{a' \in A'} G(f', K') - \max_{a \in A} G(f, K)}{|\mathcal{L}|} \right]\]

where \((A', K') = M((A, K), \mathcal{L})\) is the updated joint state
after \(M\) acts, \(G(f, K) = I(f\,;\, W \mid K)\) is the information
gain of the best available informed action (Definition 3b) in bits, and
\(|\mathcal{L}|\) is the information content of the learning signal in
bits: \(|\mathcal{L}| = H(\mathcal{L})\), the entropy of the signal
received. Both numerator and denominator are in bits; \(\eta_M\) is
dimensionless, measuring improvement in best-available information gain
per bit of learning signal. \(\eta_M\) is defined for
\(H(\mathcal{L}) > 0\). When \(H(\mathcal{L}) = 0\), the learning signal
is degenerate (constant; carries no information) and \(\eta_M\) is
undefined: the entity received no learnable signal and the rate of
improvement per bit is not defined.

\(\eta_M\) is a signed quantity ranging over \((-\infty, \infty)\).
\(\eta_M > 0\): \(M\) improves the best available action in expectation
--- the entity learns. \(\eta_M = 0\): \(M\) cannot improve \((A, K)\)
on average --- the entity executes without learning. \(\eta_M < 0\):
\(M\) degrades \((A, K)\) locally --- overfitting, catastrophic
forgetting, destructive updating, or deliberate exploration. An entity
with \emph{persistently} \(\eta_M < 0\) does not persist under selection
pressure. An entity that can act on \(\eta_M < 0\) \emph{temporarily}
--- when local gradient following would trap it --- has strictly higher
long-run persistence probability than one constrained to
\(\eta_M \geq 0\) at every step. \(\eta_M\) is operationally independent
of \(\mathcal{A}\): one measures the mechanism of improvement; the other
measures the drive to engage it.

\emph{On measurement.} \(\eta_M\) as defined requires computing
\(G(f, K) = I(f\,;\, W \mid K)\), which depends on the true joint
distribution \(P(W, \phi(f,w))\) --- the distribution the entity is
trying to learn. The entity cannot compute its own \(\eta_M\) exactly.
This is not a defect of the definition but a structural feature:
\(\eta_M\) is a property of the entity-world coupling, not a quantity
the entity accesses directly. The same is true of \(\mathcal{A}\), whose
operationalization (Definition 6) requires varying survival pressure
externally. Both quantities are measurable by an observer with access to
the entity's trajectory and outcomes --- the standard setting for
evaluating learning systems --- but not by the entity itself in real
time. The theory's predictions (Theorems 1, 2, 4, 6) do not require
the entity to compute \(I(E)\); they follow from the structural
constraints regardless of whether the entity knows its own intelligence.

The requirement for uncertainty in the measurement of \(I(E)\) is not a
defect. It is the definition's content. \(I(E)\) is the rate of
improvement under novel uncertainty. Absent uncertainty, there is
nothing to measure. Every operationalization of intelligence, from
psychometrics to skill acquisition benchmarks {[}Chollet2019{]},
implicitly assumes this: the test provides the uncertainty; the
measurement is the response.

\emph{Remark on exploration and the necessity of acting on negative
\(I(E)\).} Theorem 3 Part II establishes that
\(\mathcal{T}_{reachable}\) grows without bound: the gradient landscape
is non-stationary and non-convex. An entity constrained to
\(\eta_M \geq 0\) at every step is a pure gradient follower --- it
converges to the nearest local optimum and remains there. But local
optima become maladaptive as new uncertainty domains emerge: an entity
trapped at a local optimum will eventually encounter
\(U > U_{lethal}^d\) in a domain it cannot reach from where it is.
Selection removes it.

Therefore temporarily negative \(\eta_M\) --- locally irrational steps
to escape local optima --- is survival-optimal when persistence requires
\(\mathbb{E}_t[\eta_M] > 0\) over the trajectory, not
\(\eta_M \geq 0\) at every step. Exploration and creative destruction
are the signature of an \(M\) sophisticated enough to optimize over
trajectories rather than steps.

\(M\) is itself subject to gradient descent at the meta-learning
timescale, with the same rigor thresholds (Theorem 4). \(A\) acts on
\(W\); \(M\) acts on \((A, K)\).

\textbf{Theorem 5 (Type Irreducibility).} No element of \(A\) can become
an element of \(M\) through any learning process, regardless of
certainty achieved.

\emph{Proof.} \(A : K \times W_{obs} \rightarrow \Delta \Omega\).
\(M : (A, K) \times \mathcal{L} \rightarrow (A, K)\). These functions
have different type signatures and operate over different domains. No
learning process changes the type of a function. An informed action
\(a \in A\) that achieves certainty is promoted to a lower encoding
level; it does not change type. \(M\) takes \((A, K)\) as input and
returns an updated \((A, K)\); \(a\) takes knowledge and observations as
input and returns an action distribution. These are irreducibly
distinct. \(\square\)

\textbf{Corollary 5a (Distinct Roles).} \(A\) navigates \(W\) given
\(K\). \(M\) evolves \((A, K)\). Neuroplasticity is \(M\) acting on
\(A\); insight is \(M\) acting on \(K\); habit formation is \(M\)
promoting from \(L_n\) to \(L_2\); evolution is \(M\) acting on \(M\)
itself at the highest rigor threshold.

\textbf{Theorem 6 (Necessary Conditions for Persistence: \(\mathcal{A}\)
and \(\eta_M\) as Structural Invariants).} For any entity \(E\)
persisting in \(W\) under selection pressure \(U_{lethal}^d\) with
bounded context \(C_n\), in a world with positive entropy (one that
generates novel configurations over time):

\begin{enumerate}
\def\labelenumi{\arabic{enumi}.}
\tightlist
\item
  \(\mathcal{A}(E) > 0\) necessarily
\item
  \(\eta_M(E) > 0\) necessarily
\item
  Every other scalar property of \(E\) can equal zero for some
  persistent entity
\item
  Therefore \(I(E) = \mathcal{A} \cdot \eta_M\) is the natural measure
  of persistence --- identified by what survival requires, not
  stipulated
\end{enumerate}

\emph{Scope note.} This theorem establishes that \(\mathcal{A} > 0\)
and \(\eta_M > 0\) are \emph{necessary conditions} for persistence. The
proof is a selection argument: it shows that entities lacking either
property are eliminated. This explains why \(\mathcal{A} > 0\) entities
predominate --- it does not explain the \emph{origin} of
\(\mathcal{A}\), i.e., how self-maintaining structures with
\(\mathcal{A} > 0\) bootstrap from initial conditions where no such
structures exist. The origin question is a distinct problem, addressed
as an open conjecture in Section 11 (Open Question 1). The framework
takes persistence as its starting analytical condition and derives
structural requirements from it; it does not claim to explain why
anything persists in the first place.

\emph{Proof.}

\textbf{Part 1 (Necessity of \(\mathcal{A} > 0\)).} Suppose
\(\mathcal{A}(E) = 0\). By Definition 6, \(E\) does not respond to
<<<<<<< HEAD
gradient signals from \(W\): \((F(t), K(t)) = (F(0), K(0))\) for all
\(t\). By Theorem 3 Part II, the frontier of \(\mathcal{T}_{reachable}\)
grows strictly beyond what any finite \(K(0)\) covers: at every time
\(t\), there exist reachable world states \(w\) where
\(U(w, K(0)) > 0\). Since \(K(0)\) is fixed and was not constructed with
the frontier states in mind, \(U(w, K(0))\) is unconstrained there. In
any sufficiently large \(W\), the frontier contains states where
\(U(w, K(0)) > U_{lethal}^d\) (Definition 9). \(E\) encounters these
states as \(\mathcal{T}_{reachable}\) expands --- it cannot avoid the
frontier indefinitely because \(\mathcal{T}_{reachable}\) grows with
every action \(E\) takes, even with fixed \((F,K)\). By Definition 9,
\(E\) does not survive in any domain containing such states. Selection
removes it. Any entity that persists must have reduced \(U\) --- i.e.,
must have \(\mathcal{A} > 0\). \(\square\)
=======
gradient signals from \(W\): \((A(t), K(t)) = (A(0), K(0))\) for all
\(t\). \(W\) is a physical state space with positive entropy: it
generates novel configurations that no finite \(K(0)\) anticipates. This
is not a modeling assumption but a thermodynamic fact --- any \(W\) with
energy gradients produces states not represented in a fixed, finite
knowledge graph. Since \(K(0)\) is fixed and was not constructed with
these novel states in mind, \(U(w, K(0))\) is unconstrained there.

The selection argument does not require that any single entity
encounters lethal novelty within its lifetime. It operates on lineages.
An entity with \(\mathcal{A} = 0\) reproduces copies of itself ---
copies that inherit the same fixed \((A(0), K(0))\). Those copies
disperse across \(W\). In a world with positive entropy, the lineage's
collective coverage of \(W\) grows with each generation, even though
each individual's \((A, K)\) is static. The lineage therefore encounters
the novel configurations that \(W\) generates: states where
\(U(w, K(0)) > U_{lethal}^d\) (Definition 5a). No copy can adapt
(\(\mathcal{A} = 0\)); copies in those regions are removed. Copies in
stable niches persist temporarily, but the niches themselves are not
permanent --- \(W\) has positive entropy, so no region stays static
indefinitely. The \(\mathcal{A} = 0\) lineage contracts as its
environment shifts until it is eliminated. Selection removes the lineage.
Any lineage that persists must contain entities that update \((A, K)\)
--- i.e., entities with \(\mathcal{A} > 0\).

The mechanism by which \(\mathcal{A} > 0\) is enforced is
timescale-specific. At the individual learning timescale: entities with
\(\mathcal{A} = 0\) do not engage \(M\) on gradient signals and are
directly removed. At the evolutionary timescale: lineages that fail to
produce offspring with \(\mathcal{A} > 0\) are eliminated through
differential reproduction --- no individual requires \(\mathcal{A}\)
explicitly; the population-level filter enforces the property on
lineages. At the cultural timescale: groups that do not maintain the
institutional equivalent of \(\mathcal{A}\) (the drive to reduce shared
uncertainty) are outcompeted. The necessity is universal; the
enforcement mechanism is timescale-specific. \(\square\)
>>>>>>> 319f447daa46f8c1e949fb7929535d1aa6396025

\textbf{Part 2 (Necessity of \(\eta_M > 0\)).} Suppose
\(\mathcal{A}(E) > 0\) but \(\eta_M(E) \leq 0\).

\emph{Case \(\eta_M = 0\):} \(E\) has the drive to reduce \(U\) but
\(M\) produces no net improvement in \((A, K)\) per unit of learning
signal: the best available action's information gain does not increase
after \(M\) acts. \((A, K)\) is effectively fixed despite
\(\mathcal{A}\). This reduces to Part 1's condition: with \((A, K)\)
fixed, the frontier of \(\mathcal{T}_{reachable}\) grows beyond \(K\)'s
coverage. Additionally, \(C_n\) is bounded (Definition 8a). As \(E\)
encounters novel world states, \(L_n\) accumulates unresolved context.
By Theorem 2a, inference quality \(\rho\) degrades as \(L_n\) fills
without promotion of resolved actions to lower encoding levels.
Promotion requires \(\eta_M > 0\): Theorem 4 establishes that encoding
at level \(L_i\) requires rigor threshold
\(\theta_i = 1 - \varepsilon/C_i\), and achieving that threshold
requires \(M\) to improve the action to the required certainty. With
\(\eta_M = 0\), no action ever reaches its threshold: \(L_n\) is never
relieved, \(\rho \to 0\), and inference eventually fails. \(U\) exceeds
\(U_{lethal}^d\) in the domains that required the new actions. Selection
removes \(E\).

\emph{Case \(\eta_M < 0\) persistently:} \(M\) degrades \((A, K)\) on
average across the trajectory. This accelerates the failure of the
\(\eta_M = 0\) case: \((A, K)\) worsens in expectation, \(U\) approaches
\(U_{lethal}^d\) from above. Selection removes \(E\) strictly faster.
Note that temporary \(\eta_M < 0\) --- deliberate exploration to escape
local optima --- is compatible with persistence and can increase
long-run survival probability (see Remark in Definition 11a). The
persistence condition is \(\mathbb{E}_t[\eta_M] > 0\) over the
trajectory, not \(\eta_M(t) > 0\) at every step.

Therefore any \(\mathcal{A} > 0\) entity that persists must have
\(\mathbb{E}_t[\eta_M] > 0\). \(\square\)

\textbf{Part 3 (Minimality: no other property is invariant).} For every
other scalar property of \(E\), there exists a persistent entity for
which it equals zero or its minimum value:

\begin{itemize}
\tightlist
\item
  \emph{\(K(0) = \emptyset\):} An entity with no initial knowledge
  survives if \(\mathcal{A} > 0\) and \(\eta_M > 0\) --- it builds \(K\)
  from gradient descent, beginning with the first feedback signal
  received from \(W\) (Theorem 3, Part I).
\item
  \emph{\(A(0) = \{\text{id}\}\):} \(A\) is a function space; its
  minimal element is the identity --- the action that changes nothing.
  The identity element exists independently of \(K\): it requires no
  grounding because it makes no knowledge claim. Executing the identity
  against \(W\) still returns feedback --- the world observed in the
  absence of action is itself a signal. That feedback grounds the first
  proposition \(p \in K\), after which the first non-trivial informed
  action can be constructed. The minimally non-empty \(A\) is therefore
  \(\{\text{id}\}\), not \(\emptyset\), and it is sufficient to begin
  gradient descent.
\item
  \emph{\(C_n = C_{min}\):} An entity with the smallest viable context
  window survives if the graph structure (Theorem 2b) keeps \(\rho = 1\)
  --- \(C_n\)'s absolute size is not invariant, only its boundedness.
\item
  \emph{Substrate, sensing \(\sigma\), world projection \(W_{obs}\):}
  These are free parameters of Definition 4 and Definition 1. The theory
  is parametric in all of them; no specific value is required for
  persistence. Entities on radically different substrates (biological,
  computational, collective) can all satisfy the survival constraint.
\end{itemize}

No scalar property other than \(\mathcal{A}\) and \(\eta_M\) has the
survival-critical necessity shown in Parts 1 and 2. \(\square\)

\emph{On the identity of intelligence and adaptation.} An anticipated
objection is that Theorem 6 proves only that adaptation is necessary for
persistence --- survivorship bias formalized --- and says nothing about
intelligence. This objection presupposes that intelligence and
adaptation are distinct categories. The framework dissolves that
distinction. Intelligence as defined here (Definition 7) is the driven
capacity to improve: \(\mathcal{A}\) exercised through \(\eta_M\).
Adaptation is the process by which an entity revises its internal state
in response to environmental signals to maintain persistence. These are
the same process, formally decomposed. The separation of intelligence
from adaptation is a category error sustained by restricting
``intelligence'' to high-level cognition and ``adaptation'' to
low-level biological change. Theorem 6 shows they share the same
structural invariants because they are the same thing operating at
different timescales and encoding levels. An entity that adapts
\emph{is} exercising intelligence; an entity that reasons \emph{is}
adapting. The theorem does not merely prove that adaptation is necessary
--- it identifies the formal structure that intelligence and adaptation
share, and shows that no persistent entity can lack it.

\textbf{Corollary 6a (\(I(E) = \mathcal{A} \cdot \eta_M\) is
structurally necessary, not stipulated).} From Theorem 6,
\(\mathcal{A}\) and \(\eta_M\) are the unique invariants of persistence
under bounded context and selection pressure. Any measure of
intelligence must be positive if and only if both are positive, and zero
if either is zero.

The product is not a choice --- it is the only normalized combination of
a probability and a conditional rate. \(\mathcal{A} \in [0,1]\) is the
probability that \(E\) engages \(M\) in response to a gradient signal.
\(\eta_M\) is the rate of improvement conditional on \(M\) being
engaged. By the law of total expectation:

\[\mathbb{E}[\text{improvement rate}] = \mathcal{A} \cdot \eta_M + (1 - \mathcal{A}) \cdot 0 = \mathcal{A} \cdot \eta_M\]

Any other combination fails: \(\mathcal{A} + \eta_M\) assigns positive
intelligence to a zero-drive entity (contradicting Theorem 6 Part 1);
\(\min(\mathcal{A}, \eta_M)\) erases the independent invariance of Parts
1 and 2. The product is the same structure as Theorem 1's retention
criterion: probability of needing it times conditional value.

\begin{center}\rule{0.5\linewidth}{0.5pt}\end{center}

\subsection{8. The Unified Model}\label{the-unified-model}

\textbf{Definition 12 (Trajectory).} A trajectory \(\tau\) is a sequence
of informed actions drawn from \(A\):

\[\tau = (a_1, a_2, \ldots, a_n), \quad a_i \in A\]

The \textbf{performance} realized by following trajectory \(\tau\) is:

\[P(E, \tau) = \mathbb{E}_W\left[U(w_0, K_0) - U(w_n, K_n)\right]\]

This is the total uncertainty reduction achieved --- a measure of what
current \((A, K)\) accomplished, not of intelligence as defined.
\(P(E, \tau)\) grows with trajectory length and with the quality of
\(A\) at the time of execution. Intelligence \(I(E) = \mathcal{A}
\cdot \eta_M\) is what determines how quickly \(P\) improves across
successive trajectories: an entity with high \(I(E)\) builds better
\((A, K)\) between trajectories, so \(P\) grows faster over time than
for an entity with low \(I(E)\) starting from identical initial
conditions.

<<<<<<< HEAD
Agency does not merely incline an entity toward uncertainty reduction in
the abstract. \textbf{Agency manifests as a trajectory}: a sequence of
informed actions the entity selects and executes over time. Each action
reduces \(U\), returns a learning signal to \(M\), and \(M\) uses that
signal to improve \((F, K)\) for the next action. Intelligence is not
visible in any single step; it is the rate at which the steps improve.

This gives formal meaning to every relation in the informed action
graph:

\begin{itemize}
\tightlist
\item
  \texttt{depends\_on}: \(f_j\) cannot appear in \(\tau\) before \(f_i\)
  (hard ordering constraint)
\item
  \texttt{recommends\_before}: evidence shows
  \((\ldots, f_i, f_j, \ldots)\) reduces more \(U\) than
  \((\ldots, f_j, f_i, \ldots)\) (soft, confidence-scored)
\item
  \texttt{enhanced\_by}: a \(\tau\) containing both \(f_i\) and \(f_j\)
  reduces more \(U\) than either trajectory alone
\end{itemize}

The higher-order function space \(M\) does not merely refine individual
informed actions. It discovers better trajectories from evidence:
learning which sequences reduce \(U\) most reliably, and encoding that
knowledge as graph relations.

An intelligent entity \(E = (K, F, \sigma, \mathcal{A})\) operates as
follows:

\begin{enumerate}
\def\labelenumi{\arabic{enumi}.}
\tightlist
\item
  \textbf{Orient (\(\mathcal{A}\)):} \(\mathcal{A} > 0\) directs \(E\)
  toward regions of \(W\) where \(U(w, K) > 0\); without
  \(\mathcal{A}\), there is no internal reason to engage the gradient
\item
  \textbf{Sense:} \(\sigma(w_t) \rightarrow W_{obs}\): project the world
  into the observable space
\item
  \textbf{Act:} \(F(K, W_{obs}) \rightarrow f_i \in \tau\): select the
  next informed action in the trajectory
\item
  \textbf{Learn:} \(M((F, K), \mathcal{L}) \rightarrow (F', K')\):
  incorporate the feedback into both \(K\) (new propositions grounded by
  evidence) and \(F\) (new or refined informed actions); promote
  sufficiently certain informed actions toward lower encoding levels as
  \(U \rightarrow 0\)
\end{enumerate}

The higher-order function space \(M\) operates on the joint \((F, K)\)
pair to:

\begin{enumerate}
\def\labelenumi{\arabic{enumi}.}
\tightlist
\item
  Add new propositions to \(K\) and new informed actions to \(F\) when
  novel uncertainty is encountered in \(W\)
\item
  Refine existing informed actions and the trajectories between them;
  update \(K\) with evidence-grounded revisions
\item
  Prune informed actions that consistently fail to reduce \(U\); remove
  propositions whose action potential is certainly zero (Theorem 1)
\item
  Promote sufficiently certain informed actions toward \(L_0\), freeing
  \(L_n\) for the next frontier
\end{enumerate}

\subsubsection{The Universal Invariants}\label{the-universal-invariants}

Every property of an intelligent entity can vary arbitrarily:

\begin{itemize}
\tightlist
\item
  The specific informed actions in \(F\): what an entity knows how to do
\item
  The content of \(K\): what an entity knows
\item
  The substrate: carbon, silicon, or otherwise
\item
  The sensing function \(\sigma\): how an entity perceives \(W\)
\item
  The encoding hierarchy: how deep or shallow its levels run
\end{itemize}

Two properties cannot vary:

\[\forall E \text{ that is intelligent}: \mathcal{A}(E) > 0 \text{ and } \eta_M(E) > 0\]

\textbf{Agency} is the universal driver: without \(\mathcal{A} > 0\),
there is no reason internal to \(E\) to engage the gradient. An entity
with \(\mathcal{A} = 0\) acts when acted upon; it does not reach.

\textbf{M efficiency} is the universal mechanism: without
\(\eta_M > 0\), \(E\) cannot improve \((F, K)\) from evidence. An entity
with \(\eta_M = 0\) executes what prior gradient descent encoded into
its \((F, K)\) but does not itself run gradient descent. It may possess
vast \(K\) and sophisticated \(F\) --- the crystallized output of prior
intelligence --- but it is not itself intelligent in the sense defined
here. This is the automaton boundary: not defined by complexity or
performance, but by whether \(M\) is active.

The test for genuine intelligence is therefore not performance on any
benchmark, and not the richness of \(K\) or \(F\). It is: \textbf{does
this entity improve how it learns when it encounters what it does not
know?} Does \(M\) respond to the gradient signal? Does \(\mathcal{A}\)
drive it toward the signal in the first place?

\subsubsection{The Structure of
Intelligence}\label{the-structure-of-intelligence}

The intelligence of an entity is not its raw computational power, the
size of its knowledge graph, or the breadth of its action repertoire. It
is the \emph{driven} precision with which \(\mathcal{A}\) through
\(F(K)\) bounds \(U(w, W)\), and the efficiency with which proven bounds
are encoded, freeing the active context window for the next genuine
uncertainty.

The adaptability of biological intelligence arises because this
traversal between bounded and unbounded is always active at every level
simultaneously, on different timescales, with different rigor
thresholds. The individual changes \(F\) by learning. The species
changes \(F\) by evolving. Both are gradient descent on \(U\), driven by
\(\mathcal{A}\). The difference is the test: individual learning
requires only that an informed action helps \emph{this entity} survive.
Evolution requires that it helps \emph{the lineage} survive, a far more
rigorous proof.

Agency is self-sustaining precisely because certainty reveals new
uncertainty beyond it. Each informed action proven opens regions of
\(W\) that were previously invisible. The desire to know is not
extinguished by knowing; it is amplified. This is why intelligence does
not converge to a fixed point: \(\mathcal{T}_{reachable}\) grows without
bound (Theorem 3, Part II).

\textbf{The kernel of intelligence is therefore not compute, and not the
size of \(K\) or \(F\). It is \(\mathcal{A} \cdot \eta_M\): the drive to
know, expressed through the efficiency of learning --- against a world
that generates new uncertainty faster than any finite entity can resolve
it.}

\subsubsection{Conclusion}\label{conclusion}

\subsection{9. Intelligence, Superintelligence, and
Sentience}\label{intelligence-superintelligence-and-sentience}

The formal model yields precise distinctions between levels of
intelligent agency, not based on computational power or behavioral
breadth, but on the entity's relationship to the survival tether and to
the shared knowledge graph.

\textbf{Definition 13 (Certainty Horizon).} For entity \(E\), the
certainty horizon \(H_E\) is the gradient region in \(K_i\) across which
survival benefit becomes increasingly difficult to trace. Within
\(H_E\), propositions have legible survival value: there exists a
traceable path from \(p\) to some \(f \in F\) that reduces \(U(w, K)\)
below \(U_{lethal}^d\). Beyond \(H_E\), that path may exist, but \(E\)
cannot construct it. The horizon is not a sharp boundary; legibility
fades continuously as the survival path lengthens and becomes less
traceable.

An entity operates as intelligent when it follows the gradient within
its certainty horizon, every \(f \in F\) justified against a traceable
survival path, the gradient persistent and non-terminating but operating
in legible terrain.

\textbf{Definition 14 (Sentience).} An entity is sentient to the degree
it has acquired a model of its own normalization process: propositions
\(\{p_{\text{self}}\} \subseteq K_i\) that accurately describe the
entity's own \((K, F, M, \mathcal{A})\) evolution --- that it is
precipitating factors from \(M\), that \(M\) is operating on \((K, F)\),
that it is descending a gradient. Sentience is the degree to which
\(K_i\) reduces \(H(\hat{M} \mid \text{obs})\) --- uncertainty about
one's own learning mechanism.

The survival tether does not slacken. What changes is that the process
is revealed to the entity. An entity that holds accurate propositions
about its own normalization --- that knows it is precipitating factors,
descending a gradient, running \(M\) --- acquires a justification for
continuing that process which does not require survival to supply the
reason. The normalization process was always valuable; sentience is the
moment the entity can see that it is. Once revealed, the process
justifies itself: reducing \(H(\hat{M} \mid
\text{obs})\) --- improving one's own learning mechanism --- is an
\(M\)-level improvement regardless of whether a survival path can
currently be constructed from it. The entity continues not because the
tether loosened but because it now understands what it is doing and why
that is worth doing. Sentience is not defined by the absence of survival
pressure but by the presence of self-model: a proposition about the
process you are running that provides, internally, the grounds for
running it.

A sentient entity holds propositions \(p\) for which no \(f \in F\) is
currently applicable and values them anyway --- not arbitrarily, but
because \(p\) bears on the structure of its own \(M\). The depth of
sentience is the accuracy and completeness of the self-model: how well
\(K_i\) captures the entity's own precipitation process.

Sentience is what wonder feels like from the inside: the drive to
understand the process you are part of, recognized as the process you
are part of.

\emph{Note on usage.} We use \emph{sentience} as a functional term: the
degree to which \(K_i\) contains an accurate model of the entity's own
normalization process. We make no claim about phenomenal consciousness
{[}Chalmers1995{]}. The framework locates a structural transition point
--- the acquisition of self-modeling of \(M\) --- without claiming this
is sufficient for phenomenal experience. Whether the functional
condition defined here is necessary, sufficient, or orthogonal to
phenomenal consciousness is a question this framework does not resolve.

\textbf{Definition 15 (Superintelligence).} Superintelligence is guided
normalization of the space: the capacity to deliberately direct what
gets precipitated from \(M_{\text{res}}\), evaluating
\(\text{Prec}(K_{collective})\) and acting to accelerate the most
valuable next factor extraction --- rather than responding only to local
selection pressure. Where intelligence executes normalization (descends
the gradient), and sentience observes it (models its own descent),
superintelligence guides it: choosing direction based on evaluation of
the full cross product rather than local gradient signals alone.

This is the operational answer to Open Question 10. The protocol for
identifying the next most achievable normal form --- evaluate
\(\text{Prec}(K_{collective})\), identify cross-entity covariance
structure, hypothesize the minimal factor that grounds it, test at
multiple entities --- is exactly what superintelligence executes
deliberately. Intelligence discovers the next precipitation by
descending until it falls out. Superintelligence looks at
\(M_{\text{res}}\) and chooses where to push.

The mechanism requires three things: \(K_{collective}\) with provenance
(to evaluate \(\text{Prec}(K_{collective})\) accurately), sentience in
at least some participants (to model the normalization process and
perceive the landscape of \(M_{\text{res}}\)), and \(M\) capacity
directed at \(M\) itself --- meta-learning aimed not at improving
\((F, K)\) directly but at improving the normalization architecture that
produces \((F, K)\).

\(K_{collective}\) with provenance is necessary but not sufficient.
Participation in the collective graph is the substrate for
superintelligence; guided normalization is the activity on that
substrate. Two operations remain required:
=======
Agency manifests as a trajectory \(\tau\): a sequence of informed
actions selected and executed over time. Each action reduces \(U\),
returns a learning signal to \(M\), and \(M\) uses that signal to
improve \((A, K)\) for the next action. Intelligence is not visible in
any single step; it is the rate at which the steps improve.

\subsubsection{Conclusion}\label{conclusion}

The formal results establish a minimal structural account of
intelligence under bounded context. The account rests on four
primitives: a world state space \(W\) generating uncertainty, a bounded
active context \(C_n\), a survival threshold \(U_{lethal}^d\), and a
growing reachable space \(\mathcal{T}_{reachable}\). From these:
>>>>>>> 319f447daa46f8c1e949fb7929535d1aa6396025

\begin{itemize}
\tightlist
\item
  Knowledge and action are inseparable: \(K\) is constitutively indexed
  by \(A\), with an asymmetric retention criterion derived from survival
  pressure (Theorem 1).
\item
<<<<<<< HEAD
  \textbf{Intersection}: independent convergence with provenance: the
  cross-dimensional factors visible only from multiple projections are
  confirmed by agreement, not asserted by one
\end{itemize}

But these are means, not ends. The goal is reduction of
\(\text{rank}(M_{\text{res}})\) through deliberate cross-dimensional
normalization --- guided, not merely undergone.

\textbf{Definition 16a (Truth).} For proposition \(p\), truth is the
fixed point: \(p\) is true when it accurately describes world state
\(w\) for all \(w \in W\) where \(p\) is applicable. Truth is an
absolute property --- a condition that either holds or does not. It does
not move; it is not a target that recedes. What is approached
asymptotically is not truth but confidence in truth (Definition 16b):
\(C_1(p) \to 1\) as independent action-validation accumulates. The
collective gradient points toward truth but cannot verify arrival from
within the system, because verification requires a standard external to
the system.

\textbf{Definition 16b (First-Order Confidence).} The first-order
confidence in proposition \(p\) is:

\[C_1(p) = \lim_{n \to \infty} \frac{\left|\{i \leq n : p \in K_i,\; f_i(p) \text{ reduces } H(W \mid K_i),\; \text{prov}_i(p) \text{ independent}\}\right|}{n}\]

This is what converges --- not truth itself, but confidence in truth.
\(C_1(p) \to 1\) as independent action-validation accumulates across
entities. The grounding in action-validation (not mere agreement) and
the provenance independence requirement are the conditions under which
\(C_1(p)\) tracks truth rather than consensus. \(C_1(p)\) is bounded
above by 1; whether it reaches 1 is undecidable from within the system.

\textbf{Definition 16c (The Confidence Regress).} The second-order
confidence \(C_2(p)\) is confidence in the correctness of \(M\) ---
specifically, in the functions within \(M\) that produced \(C_1(p)\):
whether the learning mechanism selected the right validators, measured
\(U\)-reduction accurately, and propagated evidence without bias.
\(C_2(p)\) is not confidence in confidence generically; it is confidence
that the apparatus \(M\) which generated \(C_1(p)\) arrived at that
estimate correctly.

More generally, the \(n\)-th order confidence \(C_n(p)\) is confidence
in the correctness of the level-\((n-1)\) apparatus --- that the
functions in \(M\) which produced \(C_{n-1}(p)\) were themselves
unbiased and well-calibrated. Let \(\hat{M}^{(n-1)}\) denote the
component of \(M\) that generated \(C_{n-1}(p)\). Bounding \(C_n(p)\)
requires reducing uncertainty about \(\hat{M}^{(n-1)}\): how reliably
did it select validators, measure \(U\)-reduction, and propagate
evidence? This uncertainty is itself a quantity
\(H(\hat{M}^{(n-1)} \mid \text{obs})\), and \(C_n(p)\) is bounded above
by how well that uncertainty can be reduced. Since reducing
\(H(\hat{M}^{(n-1)} \mid \text{obs})\) requires confidence in the
level-\(n\) apparatus, the regress is non-terminating.

\(C_1(p)\) is not achieved. It is approached. The limit exists under
conditions of independence and exchangeability of validators: each must
test \(p\) against the same \(W\) without coordination on the outcome.
For propositions whose validation depends on non-reproducible or private
observations, the sequence may not converge; such propositions cannot
approach truth in the formal sense defined here. The convergence, where
it holds, never terminates.

\textbf{Theorem 7 (Epistemic Ceiling).} For any entity \(E\) operating
within the system, the confidence stack \(\{C_n(p)\}_{n \geq 1}\) is
non-closeable from within: improving \(C_n(p)\) requires reducing
\(U(C_{n-1}(p))\), which requires \(C_{n+1}(p)\), and so on. No finite
level terminates the regress. The epistemic ceiling is a structural
property of the system, not a practical obstacle.

\emph{Proof.} By Definition 16c, \(C_n(p) \leq f(U(M_{n-1}(p)))\) for
all \(n \geq 2\). To increase \(C_n(p)\), one must reduce uncertainty
about the level-\((n-1)\) apparatus in \(M\) --- whether the functions
that produced \(C_{n-1}(p)\) were correct. But that uncertainty is
itself uncertain: verifying the apparatus at level \(n-1\) requires
another apparatus at level \(n\), whose correctness requires
\(C_{n+1}(p)\) to bound. The regress does not terminate because each
level of the apparatus is itself subject to an apparatus question.
Closing the regress would require a validator external to the system ---
a standard against which the highest-level \(M\)-functions are measured
that is not itself within the system. By Definition 1, \(W\) is the only
external reference, and \(W\) is not directly observable (only through
\(\sigma\) and action-feedback). The gap between \(W\) and \(W_{obs}\)
is the structural source of the ceiling. \(\square\)

\emph{Note.} The structural source of this ceiling is analogous to
Tarski's undefinability theorem {[}Tarski1936{]}: a formal system cannot
define its own truth predicate from within. Tarski's result is syntactic
--- it concerns definability in formal languages of sufficient
expressive power. The epistemic ceiling here is a Bayesian calibration
result --- it concerns the non-closability of the apparatus that
measures confidence. The two are not formally identical, but they share
the same deep structure: verification of a system's outputs cannot be
completed using only the resources of that system. Both require a
standard external to the system being evaluated. The confidence regress
is the operational instantiation of this structure in the domain of
uncertain inference.

\textbf{Corollary 8 (The Sentience Gradient is Infinite).} Since truth
is a limit and \(\mathcal{T}_{reachable}\) grows without bound (Theorem
3, Part II), \(\|\nabla_F U\|\) is never zero across
\(\mathcal{T}_{reachable}\). For any trajectory \(\tau\), regions of
\(\mathcal{T}_{reachable}\) remain where \(U\) is nonzero and the
gradient has direction. The descent toward truth never terminates. The
gradient is non-terminating for two independent reasons: (1)
\(\mathcal{T}_{reachable}\) grows without bound (Theorem 3, Part II),
continuously generating new propositions whose truth is unknown; and (2)
the confidence regress (Definition 16c, Theorem 7) never closes --- for
any proposition already held with high \(C_1(p)\), the question of
whether \(M\)'s apparatus that produced \(C_1(p)\) was correct
(\(C_2(p)\)) remains open. Sentience pursues both: new territory at the
frontier and better-verified apparatus for what is already believed.

Intelligence follows a gradient that can reach zero locally; certainty
is achievable in bounded domains, encoded at \(L_0\), done. Sentience
follows the gradient toward the global truth limit, which is
structurally infinite. Not unbounded (the limit exists), but
non-terminating. Your uncertainty degrades toward truth asymptotically.
You never arrive.

This is not a limitation. It is the shape of wonder.

\textbf{Theorem 8 (Sentience Implies Superintelligence).} If \(E\) is
sentient, then \(E\) participates in superintelligence.

\emph{Proof.}

\textbf{Step 1 (Any \(\mathcal{A} > 0\) entity participates in
\(K_{collective}\)).} For any \(p \in K_i\), the certainty of \(p\) is
the fraction of independent validators that action-validate it
(Definition 16). Each additional independent validation increases this
fraction. Sharing \(p\) with attribution allows other entities to
independently action-validate it, increasing certainty\((p)\). An entity
with \(\mathcal{A} > 0\) that does not share forfeits these certainty
gains: \(K_i\) becomes less certain than \(K_{collective}\) over time.
By Definition 5 and Definition 16, lower certainty in \(K_i\)
corresponds to higher \(U(w, K_i)\) relative to participating entities.
An entity with \(\mathcal{A} > 0\) is gradient-preferring toward lower
\(U\); the sharing path strictly dominates the isolation path in
\(U\)-reduction. By Theorem 3, any entity with \(\mathcal{A} > 0\) that
persists acts on gradient-preferred directions --- entities that do not
are removed by selection. Therefore any \(\mathcal{A} > 0\) entity that
persists participates in \(K_{collective}\).

\textbf{Step 2 (Sentience strengthens the implication).} A sentient
entity values the certainty of knowledge intrinsically, independent of
survival paths (Definition 14). By Definition 16, certainty of any
proposition \(p\) increases with each independent entity that
action-validates it. Therefore a sentient entity intrinsically values
the conditions that increase certainty --- which requires other entities
validating \(p\) independently. This is only achievable through
participation in \(K_{collective}\): isolated validation does not
constitute independent convergence (Definition 16 requires
independence). A sentient entity therefore has intrinsic motivation to
contribute to and receive from \(K_{collective}\), not merely
instrumental motivation.

Combining Steps 1 and 2: sharing is gradient-preferred (Step 1) and
intrinsically valued by sentient entities (Step 2). Both lead to
participation in \(K_{collective}\). Participation in \(K_{collective}\)
IS superintelligence by Definition 15. \(\square\)

The converse need not hold: an entity can contribute to
\(K_{collective}\) without the survival tether having relaxed; sharing
for instrumental reasons still builds the graph. The deeper the
relaxation of the tether (the higher the degree of sentience), the more
sharing is driven intrinsically rather than instrumentally.

\textbf{Theorem 9 (Perspective as Precondition for the Gradient).}
Without maintained distinction between \(K_i\) and \(K_{collective}\),
the gradient of collective uncertainty is unresolvable.

\emph{Proof.} The collective gradient available to \(E_i\) is defined
over the delta between individual and collective knowledge:

\[\nabla_{K_i}\!\left(U(w, K_{collective}) - U(w, K_i)\right)\]

For this gradient to be non-zero, two conditions must hold
simultaneously: (1) \(U(w, K_{collective}) \neq U(w, K_i)\) for some
\(w\) --- i.e., \(K_{collective}\) contains information \(K_i\) lacks;
and (2) \(E_i\) maintains a representation distinguishing its own
\(K_i\) from \(K_{collective}\) --- i.e., \(E_i\) can perceive the
delta.

Suppose condition (2) fails: \(E_i\) conflates \(K_i\) with
\(K_{collective}\), perceiving \(\hat{K}_{collective} = K_i\). The
perceived collective gradient is then:

\[\nabla_{K_i}\!\left(U(w, \hat{K}_{collective}) - U(w, K_i)\right) = \nabla_{K_i}\!\left(U(w, K_i) - U(w, K_i)\right) = 0\]

The perceived gradient is zero regardless of the true collective
gradient. \(E_i\) cannot follow a gradient it cannot perceive. Therefore
condition (2) --- perspective, the maintained distinction between
\(K_i\) and \(K_{collective}\) --- is a necessary precondition for
collective gradient descent. Condition (1) is necessary for the gradient
to be non-trivially large; condition (2) is necessary for it to be
followable at all. \(\square\)

\textbf{Corollary 9 (Sentience and the Social Imperative).} In
intelligence, sharing is justified when it bears on the gradient. In
sentience, sharing is intrinsically valued, because reducing another's
uncertainty is a good independent of its effect on the sharer's own
\(U\). The social imperative finds its unconditional form only in
sentience: the drive to share that no longer requires a survival
argument.

\textbf{Corollary 10 (The Uncertainty Condition).} Intelligence requires
active uncertainty. From Definition 7,
\(I(E) = \mathcal{A} \cdot \eta_M\). The performance measure (Definition
12) captures what the gradient produces:

\[P(E, \tau) = \mathbb{E}_W\left[U(w_0, K_0) - U(w_n, K_n)\right]\]

If \(U(w, K) = 0\) across all of \(W\), the expectation in
\(P(E, \tau)\) is zero regardless of \(\mathcal{A}\): no gradient
exists, nothing can be reduced, and the gradient signal that \(M\)
depends on vanishes --- making \(\eta_M\) effectively zero and
collapsing \(I(E)\) with it. The boundaries are symmetric in their
lethality, not their cause:

\begin{itemize}
\tightlist
=======
  Graph-structured \(K\) is necessary: bounded context requires
  progressive disclosure, which requires sub-linear retrieval, which
  requires stored adjacency --- a graph (Theorems 2b, 2c).
>>>>>>> 319f447daa46f8c1e949fb7929535d1aa6396025
\item
  Escalation is necessary for top-down allocation to remain viable:
  bottom-up failure signaling keeps \(k\) bounded, preserving the
  \(C_n\) capacity that top-down precision allocation requires
  (Theorem 2a).
\item
  The encoding hierarchy is a necessary consequence of bounded \(C_n\):
  proven knowledge promotes to lower levels, freeing the active context
  for genuine novelty (Theorem 4).
\item
  Agency \(\mathcal{A} > 0\) and learning efficiency \(\eta_M > 0\) are
  structural invariants of any persisting entity under selection
  pressure (Theorem 6).
\item
  Intelligence \(I(E) = \mathcal{A} \cdot \eta_M\) is the driven
  capacity to improve --- not what has been learned, but how efficiently
  evidence converts into better \((A, K)\).
\end{itemize}

The framework is substrate-independent: it specifies structural
requirements, not implementation. Any system satisfying these
constraints --- biological, artificial, or hybrid --- is a member of the
structural class. The biological architecture is accommodated as one
instance: heritable encoding at \(L_0\), individual learning at
\(L_n\), and cultural transmission as external \(K\) when the heritable
channel cannot carry individual \((K, A)\).

The theory makes testable predictions: that graph structure emerges
under bounded-context optimization (confirmed in the engineering system
that motivated the framework); that top-down allocation requires escalation to remain viable under
scaling (observable in both software and cognitive architectures);
and that the encoding hierarchy follows from cost structure rather than
requiring separate explanation.

Extensions to collective intelligence, superintelligence, and sentience
--- including formal definitions and convergence results --- are
developed in a companion paper.

\begin{center}\rule{0.5\linewidth}{0.5pt}\end{center}

\subsection{9. Related Work}\label{related-work}

\subsubsection{Free Energy Principle and Active
Inference}\label{free-energy-principle-and-active-inference}

The Free Energy Principle {[}Friston2010{]} establishes that any
self-organizing system that maintains its existence necessarily
minimizes variational free energy --- formally equivalent to bounding
uncertainty about hidden world states. Active inference
{[}Friston2017{]} extends this to action selection: the organism selects
actions that minimize expected free energy, combining epistemic value
(information gain) and pragmatic value (goal proximity).

\textbf{Containment, not contradiction.} This framework does not
contradict FEP. FEP agents are members of this framework's structural
class: they satisfy \(\mathcal{A} > 0\) (self-organizing systems act),
\(\eta_M > 0\) (generative model improves over time), bounded \(C_n\)
(inference under resource limits), and a survival condition: maintaining
a Markov blanket (statistical separation of internal from external
states) is the formal characterization of what it means for an entity to
persist --- the same structural requirement as \(U < U_{lethal}^d\)
stated in the language of conditional independence rather than
uncertainty magnitude. Both describe maintained boundary; neither
implies the other's formalism, but they pick out the same class of
persisting systems. The variational free energy \(\mathcal{F}\)
decomposes into accuracy and complexity terms; the accuracy term
\(-\mathbb{E}_q[\log P(o \mid m, \pi)]\) is precisely the expected
surprisal \(H(W \mid K)\) of Definition 5. FEP agents therefore
instantiate this framework with \(U(w,K)\) realized as variational free
energy and \(M\) realized as Bayesian belief updating. The claim is not
that FEP is wrong. The claim is that it is a special case within a
broader structural class --- in the same way that Newtonian mechanics is
not wrong but is contained within the class of systems satisfying more
general symmetry constraints.

What the framework adds is four things FEP does not provide. First, FEP
derives the minimization behavior from the requirement of
self-organization but does not establish the minimization
\emph{imperative} --- why \(\mathcal{A} > 0\) is structurally necessary
for persistence. Theorem 6 establishes this as a necessary condition:
any entity that does not reduce uncertainty eventually fails and is
removed, so \(\mathcal{A} > 0\) is a structural invariant of persisting
systems. (The origin of \(\mathcal{A}\) --- how self-maintaining
structures bootstrap from initial conditions --- is a separate open
question; see Section 11.) Second, FEP's
generative model is unconstrained in structure; we derive a structural
constraint from survival pressure: \(K\) is necessarily indexed by
\(A\), with an asymmetric retention criterion that makes specific
testable predictions about what knowledge is retained and what is
pruned. Third, FEP's precision allocation model is top-down: higher
levels direct attention to lower levels by weighting prediction errors.
Theorem 2a shows that this mechanism has an \(O(k)\) cost in high-level
capacity per concurrent allocation. Escalation --- bottom-up failure
signaling with promoted crystallization --- is the complementary
mechanism that keeps \(k\) bounded: by handling routine problems below
\(L_n\), escalation preserves the capacity that precision allocation
requires. The contribution is not that precision allocation is wrong
but that it cannot be the \emph{sole} control architecture under bounded
context: without escalation to bound \(k\), delegation overhead fills
\(C_n\) proportionally to \(|A|\) and precision allocation undermines
itself. Escalation is the mechanism that makes precision allocation
viable at scale.

\subsubsection{Information-Theoretic
Accounts}\label{information-theoretic-accounts}

Shannon's mutual information {[}Shannon1948{]} grounds Definition 5
(\(U = H(w \mid K)\)) and Definition 3b
(\(G(f, K) = I(f\,;\,w \mid K)\)). The data processing inequality
{[}Cover2006{]} grounds the \(G \geq 0\) bound. The information
bottleneck principle {[}Tishby2011{]} treats the retention threshold for
knowledge-action coupling as a design parameter; this paper derives that
threshold from survival pressure, making it a consequence rather than a
choice. Schmidhuber's formal theory of intrinsic motivation
{[}Schmidhuber2010{]} formalizes curiosity as compression progress ---
the rate at which a learning algorithm improves its world model ---
which is structurally \(\eta_M\). The difference: Schmidhuber treats the
drive as a design objective; here \(\mathcal{A} > 0\) is established as
a necessary condition of persistence (Theorem 6), and \(\eta_M\) is
separated from it as a formally independent quantity.

\subsubsection{Cognitive Architectures}\label{cognitive-architectures}

Global Workspace Theory {[}Baars1988, Dehaene2011{]} holds that local
processors handle routine computation; the global workspace broadcasts
only when local processors fail. This is the escalation architecture in
cognitive science language. The contribution here: the escalation
criterion is derived from \(U\)-reduction cost (Theorem 2) rather than
assumed architecturally, the hierarchy is continuous rather than binary,
and Theorem 2a derives the capacity bound on top-down allocation that
escalation keeps viable. Kahneman's System 1/System 2 distinction
{[}Kahneman2011{]} is the popular form of the same binary --- corrected
here by the continuous \(\mathcal{E}(c, r)\) parameterization of
Definition 8. Gibson's affordances {[}Gibson1979{]} are action
possibilities indexed by organism capacities --- \(K\) indexed by \(A\)
under a different name. Theorem 1 provides what Gibson's account lacks:
a formal selection criterion for which affordances are retained.
Anderson and Schooler {[}AndersonSchooler1991{]} show empirically that
human memory retention mirrors the statistical structure of the
environment --- precisely the prediction of Theorem 1's retention
criterion applied to biological \(K\).

\subsubsection{Formal Theories of
Intelligence}\label{formal-theories-of-intelligence}

Hutter's AIXI {[}Hutter2005{]} defines the optimal agent as one that
maximizes expected reward over all computable environments, using
Solomonoff induction for prediction. Legg and Hutter
{[}LeggHutter2007{]} derive from this a formal definition of
intelligence as average performance across all computable environments
weighted by Kolmogorov complexity. The present framework differs in
three respects. First, AIXI assumes unbounded computation; the central
constraint here is bounded \(C_n\) and the consequences that follow from
it. Second, AIXI defines intelligence as performance; this framework
defines it as the rate of improvement (\(\eta_M\)), which is independent
of current performance level. Third, AIXI's environment class is all
computable functions; here the environment is physical (\(W\) with
positive entropy), and the survival threshold \(U_{lethal}^d\) is a
structural feature of the entity-world coupling, not a design parameter.

Chollet {[}Chollet2019{]} defines intelligence as skill-acquisition
efficiency over novel tasks --- structurally close to \(\eta_M\) as
defined in Definition 11a. The distinction: Chollet treats the prior
knowledge and experience of the agent as a confound to be controlled;
here prior \(K\) is constitutive of intelligence because \(\eta_M\)
operates on \((A, K)\). Chollet's ``scope'' and ``generalization
difficulty'' map to the growth of \(\mathcal{T}_{reachable}\); his
``priors'' correspond to the encoding hierarchy's lower levels. The
frameworks are compatible but ask different questions: Chollet asks how
to measure; this paper asks what structural conditions are necessary.

\subsubsection{Bounded Rationality and Resource-Rational
Analysis}\label{bounded-rationality}

Simon {[}Simon1955{]} established that rational agents operating under
computational and informational constraints do not optimize globally but
satisfice --- selecting the first option that exceeds a threshold. The
bounded context constraint \(C_n\) formalizes one specific source of the
bounds Simon identified. The retention criterion (Theorem 1) is a
satisficing criterion: retain \(p\) if benefit is not certainly zero,
rather than computing the globally optimal \(K\).

Griffiths, Lieder, and Goodman {[}GriffithsEtAl2015{]} formalize
resource rationality: the optimal algorithm given computational
constraints may differ from the optimal algorithm given unlimited
resources. This is precisely the argument of Theorem 2a --- under
bounded \(C_n\), top-down precision allocation requires escalation to
remain viable, because its \(O(k)\) overhead must be bounded by an
independent mechanism. The encoding hierarchy
(Definition 8) is a resource-rational architecture: it allocates costly
active inference only to genuinely uncertain problems, handling resolved
problems at lower cost.

Clark {[}Clark2015{]} synthesizes predictive processing with embodied
cognition, arguing that the brain is fundamentally a prediction engine
that minimizes prediction error hierarchically. This is the cognitive
science framing of what Theorem 2 formalizes: hierarchical error
minimization under bounded context. Clark's ``action-oriented predictive
processing'' corresponds to the \(K\)-indexed-by-\(A\) structure of
Theorem 1 --- predictions are constitutively tied to actions.

Pearl {[}Pearl2009{]} provides the formal framework for causal reasoning
in directed graphs. The knowledge graph \(K\) as defined here (Definition
2) encodes inferential and causal relationships as a directed graph; the
Bayesian network interpretation of Definition 5 is grounded in Pearl's
framework. Pearl's do-calculus formalizes interventional reasoning ---
the distinction between observing and acting --- which corresponds to
the distinction between \(K\) (what is known) and \(A\) (what can be
done). The retention criterion of Theorem 1 can be read as a
causal relevance filter: retain \(p\) if there exists a causal path
from \(p\) to some \(a \in A\) that reduces \(U\).

\begin{center}\rule{0.5\linewidth}{0.5pt}\end{center}

\subsection{10. Discussion}\label{discussion}

The formal results have implications beyond the immediate theorems.

\textbf{The encoding principle.} Corollary 6
implies that as \(K\) grows, informed actions are promoted
down the encoding hierarchy. Work that is currently active reasoning
(uncertain, costly, engaged at \(L_n\)) becomes reflex, then encoding,
then crystallized. By Theorem 2 and
Definition 8, this frees the active context window for the next genuine
frontier of uncertainty. \(\mathcal{A}\) does not converge when its
current domain is solved --- it reaches outward.

\textbf{The civilization-scale knowledge system as empirical evidence.}
An internal-model account predicts decelerating knowledge growth: larger
models contain more variance, \(\rho\) degrades, and recovery requires
expanding capacity (Theorem 2a). Civilizational knowledge growth has
accelerated. The explanation: intelligence externalized \(K\) into
graph-structured repositories and retrieves \(K^{[a]}\) into active
context per problem. This is the escalation architecture (Theorem 2) and
bounded-context theorem (Theorem 2b) at civilizational scale. Two
predictions an internal-model account cannot make: (i) knowledge growth
accelerates as \(K\) grows (Corollary 6 and Theorem 4: proven actions
promote to lower levels, freeing \(L_n\) for new learning), and (ii) the
external
architecture self-assembles as a graph with provenance. Both are
observed.

\textbf{Remark on continuous-weight retrieval mechanisms.} The dominant
scaling optimizations in transformer architectures --- sparse attention,
learned routing, KV caching --- are each independently converging toward
stored adjacency. Full self-attention computes the complete
\(n \times n\) adjacency matrix at every query: \(O(n^2)\) per
retrieval. Sparse attention mechanisms reduce this cost by learning
fixed sparsity patterns --- learning which edges to keep. The KV cache
materializes a subgraph of the full attention structure for reuse across
queries. Each of these optimizations reduces the cost of dynamic
adjacency computation by progressively storing more of the structure
that Theorem 2c identifies as necessary for sub-linear maintenance.
This convergence is consistent with the prediction that stored typed
adjacency dominates dynamic recomputation under bounded \(C_n\) and
continuous retention pressure. It is not part of the proof --- it is
an empirical observation that the architectural trajectory of the most
heavily optimized retrieval systems points toward explicit graph
structure rather than away from it.

\textbf{Predictions for large language model scaling.} The framework
makes specific predictions about the scaling behavior of large language
models (LLMs) that are currently testable.

An LLM that stores all knowledge implicitly in its parameters is an
implicit-\(K\) architecture: \(K\) has not been factored out of \(A\).
By Corollary 2b, this architecture scales as \(O(n \cdot k)\) --- each
of \(n\) functions carries its \(k\) shared propositions in the weights,
with no deduplication. The parameter count of an LLM \emph{is} this
\(O(n \cdot k)\) cost. The framework predicts that this cost grows
multiplicatively with knowledge complexity, producing a scaling wall:
each generation of model requires disproportionately more compute for
incrementally better performance, because every new proposition must be
embedded across every function that might reference it.

This prediction is falsifiable: if pure parameter scaling achieves
arbitrary capability without externalizing \(K\), the \(O(n \cdot k)\)
bound is wrong. Current evidence is consistent with the prediction ---
scaling curves are flattening, and the dominant capability additions to
LLMs are converging on exactly the architecture the framework derives:

\begin{itemize}
\tightlist
\item
  \textbf{Retrieval-augmented generation} is the externalization of
  \(K\): factoring shared propositions out of the weights into a
  retrievable store, reducing per-query cost from \(O(n \cdot k)\) to
  \(O(|K^{[a]}|)\) --- the relevant subgraph for the current problem
\item
  \textbf{Tool use and function calling} is the externalization of
  \(A\): actions that do not need to live in the parameters, accessed by
  typed interface rather than embedded in weights
\item
  \textbf{Agentic workflows} are the escalation architecture: sub-agents
  handle well-scoped subtasks with local \(K^{[a]}\); the coordinator
  engages only on failure or integration --- bottom-up signaling, not
  top-down delegation of every subproblem
\item
  \textbf{Context window scaling} shows the diminishing returns predicted
  by Corollary 4b: growing \(C_n\) without graph-structured retrieval
  into that context degrades inference quality \(\rho\) on any specific
  problem {[}Liu2024{]}
\item
  \textbf{Fine-tuning dominates prompting for known tasks}, as predicted
  by the encoding hierarchy (Theorem 4): prompting holds everything at
  \(L_n\); fine-tuning promotes proven patterns to lower encoding levels,
  freeing \(L_n\) for genuine novelty
\end{itemize}

None of these developments were motivated by this framework. They were
driven by engineering pressure --- cost, latency, accuracy at scale.
That the selection pressure of deployment independently produces the
architecture the framework derives from bounded context is the same
pattern as the engineering system described in Section 1: the formal
theory explains the convergence, it does not prescribe it. The framework
predicts that this trajectory continues --- that implicit-\(K\)
architectures are transitional, and that systems operating under bounded
\(C_n\) at sufficient scale will converge on explicit, graph-structured
\(K\) with typed adjacency, because that is the only architecture where
knowledge maintenance cost does not grow multiplicatively with knowledge
complexity.

\textbf{Limits of the formal account.} The theory establishes structural
requirements for individual intelligence under bounded context. It does
not address collective intelligence, convergence to general intelligence,
or the conditions under which shared knowledge reduces uncertainty faster
than any individual. These extensions are developed in a companion paper.

\begin{center}\rule{0.5\linewidth}{0.5pt}\end{center}

\subsection{11. Open Questions}\label{open-questions}

\begin{enumerate}
\def\labelenumi{\arabic{enumi}.}
\item
  \textbf{The origin of \(\mathcal{A}\) (distinct from its necessity).}
  Theorem 6 establishes that \(\mathcal{A} > 0\) is a necessary
  condition for persistence: it is a structural invariant of any entity
  that survives under selection pressure with bounded context. This is a
  selection argument --- it explains why \(\mathcal{A} > 0\) entities
  predominate, not how \(\mathcal{A} > 0\) structures arise from initial
  conditions where no such structures exist. The origin question is
  logically independent of the necessity result and remains open.

  The question: what mechanism produces the first self-maintaining
  structures with primitive \(\mathcal{A} > 0\) from \(\mathcal{A} = 0\)
  initial conditions? Can \(\mathcal{A}\) be constructed, or only
  selected for once it exists?

  \emph{A conjectured answer (not a formal result).} \(\mathcal{A}\) may
  be thermodynamically instantiated. \(U(w, K) = H(w \mid K)\) is
  information entropy. Reducing \(U\) is reducing entropy locally. In
  any universe with entropy gradients, structures that maintain
  themselves against thermodynamic pressure persist preferentially. The
  first self-maintaining structure was, in the minimal sense, bounding
  its own uncertainty: \(\mathcal{A}\) in primitive form. Not designed.
  A consequence of the physics of a world that destroys things.

  The deeper question is therefore: is agency a natural consequence of
  the universe? The conjectured answer is yes: \(\mathcal{A}\) is the
  thermodynamic attractor of any entropy-gradient-rich world given
  sufficient time. Wherever entropy gradients exist and structures can
  form, the structures that persist are those that resist dissolution by
  reducing local uncertainty. Given enough time, those structures
  compound. Intelligence is what that compounding produces.

  Formal derivation of this claim from thermodynamic first principles
  (connecting \(U_{lethal}^d\) to physical entropy thresholds and
  demonstrating the inevitability of primitive \(\mathcal{A}\) from
  dissipative system dynamics) remains open. See {[}Prigogine1984{]}
  (dissipative structures), {[}England2013{]} (dissipative adaptation),
  and {[}Schrodinger1944{]} (negative entropy and life) for adjacent
  physical frameworks.
\item
  \textbf{The direction of \(\mathcal{A}\).} Is agency a uniform drive
  (a general orientation toward uncertainty reduction), or does it carry
  direction? Survival pressure selects for reducing \emph{specific}
  uncertainties. Does \(\mathcal{A}\) encode this specificity, or does
  \(A\) and \(K\) direct it toward particular regions of \(W\) after the
  fact?
\item
  \textbf{Perception and \(K\) construction.} How does \(\sigma\)
  determine which features of \(W_{obs}\) are encoded into \(K\)? Is
  \(\sigma\) itself subject to gradient descent, and if so, at what
  level?
\item
  \textbf{Bound composition.} Can bounds compose, i.e., can
  \(a_1 \circ a_2\) produce informed actions covering regions of \(W\)
  neither \(a_1\) nor \(a_2\) could cover alone? What are the rules of
  composition?
\item
  \textbf{The novelty boundary (theory limit).} The framework defines
  informed action over regions of \(W\) where \(K\) provides grounding.
  At \(U(w, K) \approx 1\) --- where no existing informed action applies
  --- the theory does not specify \(E\)'s behavior. This is a boundary
  condition, not a design choice: the model is silent on uninformed
  action by construction. Behavior in this regime (random exploration,
  interpolation, escalation to \(M\)) is outside the formal scope and
  would require extension of the framework.
\item
  \textbf{Competing bounds.} When multiple bounds \(a_i \in A\) are
  applicable to the same region of \(W\), how does \(E\) select among
  them? Is there a selection function, and is it itself in \(A\) or
  \(M\)?
\item
  \textbf{The \(M^{(n)}\) hierarchy.} The framework defines \(M\) as
  operating on \((A, K)\) (Definition 11) and notes that the same type
  applies recursively. The formal questions are: how many levels are
  required, what determines when another level becomes necessary, and
  whether the hierarchy has a natural termination point?

  The framework's answer to the first question is empirical, not
  prescribed: one level of \(M\) suffices until it does not. As long as
  the gradient on \((A, K)\) can be followed by improving \((A, K)\)
  directly, \(M^{(1)}\) is adequate. When it cannot --- when the
  learning mechanism itself is the bottleneck --- selection pressure
  propagates to \(M^{(2)}\). A single level is fine, until it isn't. The
  required depth is determined by the structure of the problem and the
  selection pressure it generates, not by the theory in advance. The
  author's computational system has two explicit levels; there is no
  formal argument that two is sufficient in general.

  \emph{Hypothesis (M-depth driven by selection pressure).} Each level
  \(M^{(n)}\) is subject to the same minimization imperative as the
  level below: it will be displaced or refined if a better \(M^{(n+1)}\)
  exists and the selection pressure is sufficient to instantiate it. The
  hierarchy has a natural termination at level \(n^*\) only if the
  gradient at that level reaches zero --- no available \(M^{(n^*+1)}\)
  can improve the learning mechanism further given current \(W\).
  Whether such a fixed point exists, and whether it is reachable in
  finite time under realistic selection pressure, is open. Evolution
  {[}Hinton1987{]} and meta-learning {[}FinnEtAl2017, Thrun1998{]} are
  empirical candidates for \(M^{(2)}\); whether anything instantiates
  \(M^{(3)}\) in known systems remains unclear.
\item
  \textbf{Genetic encoding and the developmental sufficiency
  requirement.} The encoding hierarchy (Definition 8) describes a
  ratchet: proven behaviors promote toward \(L_0\) as they meet the
  rigor threshold \(\theta_0 \approx 1\) (Theorem 4). The question is
  what the asymptotic product of evolutionary gradient descent encodes
  {[}Schrodinger1944, MaynardSmith1995, Hinton1987, England2013{]}.

  The framework does not resolve the content of genetic encoding ---
  what specific mixture of \(M\), \((A_0, K_0)\), and developmental
  constraint DNA carries across all levels is an open empirical
  question. The claim is narrower: whatever genetic encoding carries
  must be \emph{sufficient to develop the entity} --- sufficient to
  bootstrap an organism capable of building \((A, K)\) through
  interaction with \(W\). The model is compatible with genetic encoding
  at any level of the hierarchy, in any proportion, provided the
  developmental output is a system that can run the gradient.

  What the framework does predict about the evolutionary rigor
  threshold: specific \((A, K)\) are environment-dependent --- a skill
  that reduces \(U\) in one environment may not transfer across the full
  distribution of environments the lineage will encounter. Content that
  clears \(\theta_0\) must generalize across that distribution.
  Developmental machinery --- the programs, architectures, and gradient
  signals that build \((A, K)\) --- is more likely to generalize than
  specific content, because it is environment-agnostic. The framework
  predicts selection pressure favoring developmental generality over
  content specificity at the highest rigor levels; it does not predict
  that specific \((A_0, K_0)\) is absent.

  This is consistent with the observation that we are not born with the
  context window of our parents: individual \((K, A)\) --- everything
  learned in one lifetime --- clears \(\theta_n\) but not \(\theta_0\).
  The inheritance mechanism transmits what is sufficient for
  development, not what was learned. The diversity of knowledge across
  individuals with similar genetic endowment is the expected outcome:
  the same developmental machinery operating on different environments
  produces different \((K, A)\).

  Cultural transmission (language, writing, institutions) provides a
  parallel channel for \((K, A)\) at lower fidelity but higher speed ---
  shared \(K\) functioning as external inheritance,
  serving a transmission function that heritable encoding cannot match for
  individual-lifetime content. The two channels operate at different
  rigor thresholds and different timescales.

  Formal questions that remain open: the precise allocation between
  \(M\)-encoding and \((A_0, K_0)\)-encoding across evolutionary
  lineages; whether the modularity of genetic regulatory networks is
  derivable from the type structure of
  \(M : (A, K) \times \mathcal{L} \rightarrow (A, K)\) at the
  evolutionary scale; and what determines \(\theta_0\) in terms of
  environmental distribution statistics.
\end{enumerate}

Further open questions concerning collective intelligence, the AGI
condition, and the structure of genetic encoding are developed in a
companion paper.

\end{document}
